\documentclass[a4paper,book,masters,en,listoffigures,listoftables]{NMBU}

\usepackage{graphicx}
\usepackage{float}
\usepackage{amsmath,amssymb}
\usepackage{chngcntr}
\usepackage{booktabs}
\usepackage{tabularx}
\usepackage{subcaption}
\usepackage{tikz}
\usetikzlibrary{arrows.meta, positioning, shapes.geometric, fit, calc}


\addbibresource{references.bib}

\counterwithout{equation}{chapter}
\counterwithin{table}{section}
\renewcommand{\theequation}{\arabic{equation}}

\title{Automated Prediction of Mucus Plug Burden in Lung CT Scans}
\author{Taofik Muhriz}
\thesisyear{2026}
\credits{30}
\faculty{Faculty of Science and Technology}
\studyprogramme{Data Science}
\supervisor{Fadi Al Machot}

\abstract{%
Mucus plugs are clinically relevant airway findings that can be assessed from chest CT, but manual assessment requires careful review of volumetric image data. This thesis investigates whether patient-level mucus plug burden can be predicted automatically from lung CT scans using a weakly supervised, slice-based machine learning pipeline. The task is formulated as patient-level regression: the available target is one mucus-burden score per CT examination, while the model operates on selected axial slices without slice-level mucus annotations.\par
The implemented pipeline combines representative slice selection, CT input representation, CNN-based feature extraction, patient-level aggregation, and transfer-oriented evaluation. Representative slices are selected using central restriction, heuristic informativeness filtering, PCA-based descriptor reduction, and MiniBatchKMeans diversity sampling. The model input is evaluated using a two-window CT baseline and a three-channel representation with a high-pass channel. ResNet18 and DenseNet121 are compared as compact pretrained CNN backbones, and an ORB-based bag-of-visual-words branch is evaluated as a handcrafted feature-fusion component. Internal performance is assessed using five-fold patient-level cross-validation. Additional transfer experiments with MosMed are used to study source--target adaptation and domain shift.\par
The results suggest that representative slice selection performed better than uniform sampling, and that the ResNet18 three-channel model with ORB-BoVW fusion gave the strongest internal performance among the tested configurations. DenseNet121 did not consistently outperform ResNet18. In the transfer experiments, fine-tuning was more useful than zero-shot transfer. MosMed is not used as external validation for mucus-burden prediction because it has a different COVID-related severity target, it is used only for transfer and domain-shift analysis. The pipeline gives promising exploratory results, but the conclusions are limited by the small internal cohort, weak supervision, and lack of external mucus-specific validation.
}

\sammendrag{%
Slimpropper er klinisk relevante funn i luftveiene og kan vurderes fra lunge-CT, men manuell vurdering krever gjennomgang av volumetriske bildedata. Denne masteroppgaven undersøker om slimproppbelastning på pasientnivå kan predikeres automatisk fra lunge-CT ved hjelp av en svakt veiledet, snittbasert maskinlæringspipeline. Oppgaven er formulert som regresjon på pasientnivå, hvert CT-opptak har én tilhørende slimproppscore, mens modellen arbeider med utvalgte aksiale snitt uten snittvise annotasjoner av slimpropper.\par
Den implementerte pipelinen kombinerer representativt snittutvalg, CT-basert inputrepresentasjon, CNN-basert egenskapsuttrekk, aggregering til pasientnivå og transfer-orientert evaluering. Representative snitt velges ved hjelp av sentral begrensning, heuristisk filtrering av lite informative snitt, PCA-basert reduksjon av snittdeskriptorer og MiniBatchKMeans-basert diversitetsutvalg. Modellinput evalueres med en to-vindus CT-representasjon og en trekanals representasjon med en high-pass-kanal. ResNet18 og DenseNet121 sammenlignes som kompakte forhåndstrente CNN-arkitekturer, og en ORB-basert bag-of-visual-words-gren evalueres som et håndlaget egenskapstillegg til CNN-representasjonen. Intern ytelse vurderes med femfolds kryssvalidering på pasientnivå. I tillegg brukes transfer-eksperimenter med MosMed for å undersøke tilpasning mellom datasett og domeneskift.\par
Resultatene tyder på at representativt snittutvalg ga bedre resultater enn jevn sampling, og at ResNet18-modellen med trekanals input og ORB-BoVW-fusjon ga den sterkeste interne ytelsen blant de testede konfigurasjonene. DenseNet121 ga ikke en konsekvent forbedring over ResNet18. I transfer-eksperimentene var finjustering mer nyttig enn zero-shot transfer. MosMed brukes ikke som ekstern validering av slimpropprediksjon, fordi datasettet har et annet COVID-relatert alvorlighetsmål, det brukes kun til analyse av transfer og domeneskift. Samlet viser pipelinen lovende utforskende resultater, men konklusjonene begrenses av liten intern kohort, svakt veiledet læring og manglende ekstern validering med slimproppspesifikke mål.
}

\acknowledgements{%
I would first and foremost like to thank my supervisor, Fadi Al Machot, for his excellent guidance, helpful feedback, and valuable support throughout the work on this master’s thesis. I could not have done this without his help. His feedback has been very useful throughout the entire process and has contributed to further developing the thesis. I also greatly appreciate that he has taken the time to discuss ideas, answer questions, and give advice when I have been stuck.

I would also like to thank my family for their support, patience, and encouragement throughout my studies. It has meant a lot to have people around me who have supported me both during periods of progress and during periods when the work has been demanding.

In addition, I would like to thank my friends and fellow students for good conversations, motivation, and support along the way. You have made everyday student life at NMBU both easier and more rewarding. I am grateful to everyone who, in different ways, has contributed to making my years at NMBU a valuable and educational period.
}

\begin{document}

\startthesis
\pagestyle{plain}

\chapter{Introduction}

\section{Clinical Motivation}

\subsection{Why CT-Based Mucus Plug Assessment Matters}
Mucus plugs are intraluminal accumulations of abnormally retained mucus that partially or completely obstruct the airways \cite{fahy2010airwaymucus,dunican2018autopsy}. In healthy lungs, mucus contributes to airway defense and is normally cleared through coordinated mucociliary transport \cite{fahy2010airwaymucus}. In chronic airway disease altered mucus composition, impaired clearance, and inflammatory processes can transform a normally protective material into a persistent source of obstruction \cite{fahy2010airwaymucus,dunican2018autopsy}. In asthma, autopsy and imaging studies indicate that mucus plugging is not restricted to rare terminal events, but can also be present as a clinically relevant feature of chronic disease \cite{dunican2018autopsy}. Greater CT-measured mucus plug burden has been associated with worse airflow obstruction and with eosinophilic or type 2 airway inflammation, which makes mucus plugging relevant not only as an imaging finding but also as an indicator of disease burden \cite{dunican2018mucus}.


Computed tomography (CT) is relevant in this context because it provides noninvasive visualization of airway and lung structure and has long been used to study structural changes related to respiratory dysfunction \cite{dejong2005ctairways}. This is important for mucus plugging because the abnormality is distributed within the bronchial tree rather than confined to a single projection image or laboratory measurement. On CT, mucus plugs may appear as complete bronchial occlusions seen as tubular, branching, or rounded opacities that can be followed across sequential slices and related to the bronchial lumen \cite{dunican2018autopsy,dunican2018mucus}. CT therefore supports structured assessment of plugging throughout a full volumetric examination, rather than limiting interpretation to isolated image examples or local observations \cite{dejong2005ctairways,dunican2018mucus}. Raw axial CT slices from the study data are shown in Figure~\ref{fig:raw_ct_examples} to illustrate the imaging context in which mucus plug assessment is performed.

\begin{figure}[H]
    \centering
    \includegraphics[width=0.32\textwidth]{figures/rawct1.png}\hfill
    \includegraphics[width=0.32\textwidth]{figures/rawct2.png}\hfill
    \includegraphics[width=0.32\textwidth]{figures/rawct3.png}
    \caption[Representative raw chest CT slices from the study data]{Representative raw axial chest CT slices from the study data shown in lung windowing. The figure illustrates the imaging context in which mucus plug assessment is performed and emphasizes that clinically relevant information is embedded within a volumetric examination rather than a single isolated image.}
    \label{fig:raw_ct_examples}
\end{figure}

A patient-level mucus score is therefore meaningful as a compact summary of disease-relevant information across the examination. In a segment-based CT scoring approach, the lungs are assessed systematically and the presence or absence of mucus plugging is summarized into an aggregate score over the full scan \cite{dunican2018mucus}. Such a score is clinically useful because it reflects overall burden at the level at which patients are followed, compared, and interpreted in relation to lung function and inflammatory status \cite{dunican2018mucus}. At the same time, deriving this score manually requires systematic review across multiple bronchopulmonary segments and sequential CT slices, which makes the process time-consuming and dependent on careful volumetric interpretation.

The clinical motivation for this thesis follows from this setting. Existing CT-based mucus plug work shows that the target is clinically relevant, but leaves open how it can be estimated computationally under realistic data constraints. The purpose of this thesis is therefore not to replace clinical judgment with a single image-level shortcut, but to investigate whether chest CT can support patient-level scoring through a constrained and reproducible machine learning formulation.

\subsection{Problem Statement}

This thesis addresses automated patient-level prediction of mucus plug burden from chest CT. The central technical difficulty is that the clinically meaningful target is defined for the patient as a whole, whereas the available image evidence is observed at slice level. The task is therefore not to classify one image in isolation, but to infer a single patient-level score from a volumetric examination in which relevant findings may appear unevenly across slices.

This creates a weakly supervised problem setting. Weak supervision arises when the available labels are coarser than the image evidence used for learning \cite{ren2023weakly}. In the present case, the model processes slice-level image content, but the available supervision is defined only at the patient level. The task is therefore more challenging than settings in which the training labels directly identify the relevant image content, because the model must learn how local observations contribute to a patient-level target without slice-level annotations \cite{ren2023weakly,qi2021drmil}.

The problem is further constrained by data availability. In a small internal cohort, not every slice is equally informative for patient-level prediction, and model design choices become closely tied to the risk of overfitting and unstable generalization \cite{dhiman2023samplesize}. Existing clinical CT scoring literature establishes that mucus plug burden can be assessed meaningfully at patient level, while related weakly supervised CT literature provides methodological ideas for learning from coarse case-level labels. What remains less well addressed is the combination of these perspectives, a patient-level mucus score prediction from chest CT under weak supervision and limited data. This thesis is positioned within that gap.


\section{Thesis Aim and Research Questions}

\subsection{Objective}
The objective of this thesis is to investigate whether patient-level mucus plug burden can be estimated from chest CT using a reproducible machine learning pipeline that remains realistic with respect to the available supervision and dataset size. The thesis studies how representative slice selection, CT input representation, feature fusion, compact CNN backbones, and transfer-oriented evaluation can support patient-level prediction from volumetric CT. The transfer experiments use MosMed as an external CT dataset for representation transfer and domain-shift analysis, not as external validation of mucus-burden prediction, because MosMed has a different COVID-related severity target.

The aim is therefore not to propose a new clinical scoring system, nor to claim a clinically validated decision-support tool. Instead, the thesis seeks to evaluate a constrained computational approach to a clinically meaningful CT scoring problem and to clarify which design choices are useful, defensible, and informative within this setting.


\subsection{Research Questions}
The thesis is guided by the following research questions:

\begin{description}
    \item[\textbf{RQ1.}] To what extent can a slice-based weakly supervised regression pipeline predict patient-level mucus plug burden from chest CT?
    
    \item[\textbf{RQ2.}] How does representative slice selection affect patient-level prediction compared with simpler slice sampling?
    
    \item[\textbf{RQ3.}] How do alternative CT input representations and feature-fusion strategies, including CNN-based learned representations and bag-of-visual-words-style handcrafted features, compare under limited-data conditions?
    
    \item[\textbf{RQ4.}] How do compact pretrained CNN backbones compare for patient-level mucus plug burden prediction from chest CT?
    
    \item[\textbf{RQ5.}] What do transfer-learning experiments with external chest CT data and a different target reveal about domain shift, adaptation, and the limits of generalization in this small-cohort setting?
\end{description}

These research questions are formulated to be empirically answerable within the scope of the thesis. They connect the clinical motivation for patient-level mucus scoring to the methodological decisions evaluated in the later chapters.


\subsection{Scope and Delimitations}
The scope of the thesis is limited to patient-level prediction from chest CT. It does not perform segmentation of mucus plugs, localization of individual plugged airways, or detection of specific bronchopulmonary segments. It also does not include full 3D model training, large-scale training, or a clinically ready decision-support system. The contribution is a reproducible machine learning pipeline and empirical analysis of its design choices, not a validated clinical product.

The main delimitations follow directly from the data and project scope. The available labels are patient-level scores, so the thesis cannot make slice-level claims about where mucus plugs are located unless such information is introduced in future annotation work. The study is also conducted under limited-data conditions, which necessarily constrains the strength of claims about generalization beyond the present setting \cite{dhiman2023samplesize}. The experiments evaluate design choices within the implemented pipeline, they do not exhaust the full space of possible medical-imaging models or clinical predictors.


\section{Contribution Summary}

Given these boundaries, the thesis makes the following contributions:

\begin{itemize}
    \item it formulates patient-level mucus plug burden prediction from chest CT as a weakly supervised regression problem.
    
    \item it develops a reproducible slice-based prediction pipeline that separates preprocessing, representative slice selection, input representation, feature extraction, and patient-level aggregation under limited-data conditions.
    
    \item it evaluates how key design choices, including slice-selection strategy, CT input representation, feature fusion, CNN backbone choice, aggregation, and transfer-oriented evaluation, affect patient-level prediction performance.
    
    \item it provides an empirical analysis of the strengths and limitations of this formulation in relation to weak supervision, small-cohort learning, and clinically relevant patient-level scoring.
\end{itemize}

The contribution of the thesis is therefore methodological and empirical rather than clinical in the narrow sense. It investigates a comparatively underexplored combination of target, supervision setting, and computational design within a reproducible machine learning framework.


\section{Thesis Outline}
The remainder of the thesis is organized as follows. Chapter~2 reviews the literature most relevant to this work, with emphasis on CT-based mucus plug assessment, learning approaches for patient-level CT analysis, and the research gap addressed by the thesis. Chapter~3 presents the theory and problem formulation, including the patient-level prediction task, the label structure, weak supervision, slice representation, feature learning, aggregation, regression, and evaluation. Chapter~4 describes the implemented methodology, covering dataset and target definition, preprocessing, representative slice selection, input representation, model architecture, patient-level aggregation, training, transfer learning, and reproducibility. Chapter~5 presents the experimental protocol and quantitative results from the internal cross-validation and transfer-learning experiments. Chapter~6 discusses the findings in relation to the research questions, including the ablation-style interpretation of pipeline components, weak supervision, low and negative \(R^2\) values, transfer/domain-shift results, and the main limitations of the study. Finally, Chapter~7 concludes the thesis by summarizing the main contributions, limitations, future directions, and final remarks.



\chapter{Related Work}

This chapter reviews the literature most relevant to the thesis. It does not aim to survey medical imaging broadly, but focuses on two areas that frame the present work. The first is CT-based mucus plug assessment, including clinical scoring and recent automated analysis. The second is methodological work on weakly supervised CT learning, multiple-instance learning, and transfer learning in limited-data settings. The chapter ends by identifying the gap addressed by this thesis, a patient-level mucus-burden prediction from chest CT when only coarse patient-level labels are available.


\section{CT-Based Mucus Plug Assessment}

\subsection{Clinical CT Scoring of Mucus Plug Burden}
The literature most closely aligned with this thesis is the CT-based mucus-plug scoring literature in asthma. Dunican et al.\ established a segment-based visual CT score and showed that greater mucus plug burden was associated with worse airflow obstruction and eosinophilic or type 2 airway inflammation \cite{dunican2018mucus}. This work is important here because it shows that mucus plug burden can be assessed systematically from CT and interpreted at the patient level, rather than only as a localized imaging finding. The related review by Dunican, Watchorn, and Fahy places mucus plugging in a broader clinical context by linking it to airway obstruction and disease mechanisms in asthma \cite{dunican2018autopsy}. Together, these studies frame mucus plugging as a clinically relevant burden measure that can be assessed from chest CT.

For this thesis, the key point is that the target is not the presence or absence of one visible plug, but the aggregate burden of plugging across the lungs. The original clinical scoring literature therefore justifies the patient-level endpoint used in this work. What it does not address is how such a score might be estimated automatically when the image evidence is distributed across a volumetric scan and only coarse labels are available.


\subsection{Recent Clinical and Pathophysiologic Work on Mucus Plugs}
More recent clinical literature suggests that mucus plugs remain an active pathophysiologic and imaging-relevant research theme beyond the original CT scoring studies. Hsieh et al.\ report that mucus plugs correlate with small airway remodelling in asthma, which strengthens the broader interpretation of mucus plugging as a marker of meaningful airway abnormality rather than only a descriptive CT finding \cite{hsieh2025smallairwaymucus}. This is important because it reinforces the idea that mucus-plug burden reflects underlying airway pathology that is relevant to disease severity and structural change.

From the perspective of the present thesis, such work helps broaden the significance of the target. The thesis does not aim to contribute directly to airway biology or histopathologic interpretation, but it benefits from a literature that continues to show that mucus plugs are clinically meaningful and mechanistically connected to disease processes. This matters because a computational prediction problem is more compelling when the quantity being predicted is already supported by a broader clinical and pathophysiologic literature.


\subsection{Recent Computational Work on Automated Mucus Plug Analysis}
Recent work has also begun to address automated computational analysis of mucus plugs on chest CT more directly. van der Veer et al.\ study automatic AI-based quantification of airway-occlusive mucus plugs in COPD and relate that quantification to all-cause mortality \cite{vanderveer2025aicopdmucus}. This is relevant because it demonstrates that automated quantification of mucus-plug burden is not only technically feasible, but can also be tied to clinically important outcome analysis. The paper therefore strengthens the broader case that automated mucus-plug analysis can have real clinical relevance rather than being only a technical image-processing exercise.

Raut et al.\ provide a different kind of contribution by validating an AI-based automated PRAGMA and mucus plugging algorithm in pediatric cystic fibrosis \cite{raut2025pragmaai}. Here the relevance lies in showing that automated mucus-plug analysis is being pursued across disease settings and that recent work is increasingly concerned with reproducible, validated computational pipelines rather than only exploratory proof-of-concept models. Although pediatric cystic fibrosis differs substantially from the present thesis setting, the paper still supports the broader point that automated CT-based mucus-plug analysis is becoming an active area of applied research.

Pu et al.\ move even further toward automated lesion-level analysis by presenting a simulation-driven, annotation-free deep learning framework for automated detection and segmentation of airway mucus plugs on thin-slice non-contrast chest CT \cite{pu2026simulation}. Their approach is notable because it reduces dependence on voxel-wise manual training annotations by generating synthetic plug examples within segmented airways. This is methodologically interesting for the present thesis because it shows another way in which the field is trying to overcome annotation bottlenecks in mucus-plug analysis. At the same time, the task formulation remains different from that of the present work, Pu et al.\ address plug-level detection and segmentation, not patient-level prediction of mucus burden from weak labels.

Together these newer studies show that automated mucus-plug analysis from chest CT is now an active recent direction. They span several related but distinct computational goals: automated quantification, disease-specific validation, and lesion-level detection or segmentation. This is an important development in the literature because it indicates that the field is beginning to move beyond purely visual or manually scored assessments toward more reproducible computational pipelines. At the same time, the methodological diversity of these studies also highlights that there is not yet one settled computational formulation of the problem.


\subsection{Limits of Existing CT Mucus-Plug Literature for the Present Task}
Despite this progress, the scope of the existing CT mucus-plug literature still differs from the scope of the present thesis. The original clinical scoring papers are primarily interpretive, they define and apply a patient-level burden score and evaluate its physiological relevance. More recent computational work extends the literature toward automated quantification, validation, detection, and segmentation, but it still does not primarily address automated patient-level regression from chest CT under weak supervision.

This distinction matters because the present thesis does not attempt to replace the clinical literature with a new scoring definition. Rather, it takes an existing clinically meaningful burden concept as its starting point and asks whether a score of this general type can be estimated computationally from CT using a constrained machine learning pipeline. The gap is therefore not in the existence of CT-based mucus-plug scoring itself, but in the limited literature connecting that scoring logic to weakly supervised patient-level prediction.

From the perspective of this thesis, the existing literature establishes the relevance of the target but does not resolve how a predictive model should handle volumetric CT data when only coarse patient-level labels are available. Clinical scoring papers justify the endpoint, newer automated studies show that computational mucus-plug analysis is feasible, but neither group of studies directly addresses the specific setting considered here, which is a patient-level mucus-burden prediction from chest CT under weak supervision and limited data. The present work is positioned in that methodological gap.



\section{Learning Approaches Relevant to CT-Based Patient-Level Assessment}


\subsection{Weakly Supervised and Multiple-Instance CT Learning}
Weakly supervised medical-imaging studies typically address the problem that clinically meaningful labels are available at a coarser level than the image evidence itself \cite{ren2023weakly}. In such settings, the methodological challenge is not only to learn from limited supervision, but also to determine how local image evidence should be related to a case-level target when the relevance of individual slices, patches, or regions is not known in advance. This has made multiple-instance and bag-level formulations particularly relevant in medical imaging, because they provide a natural way to represent an examination as a collection of local instances while keeping prediction defined at the level of the whole case.

In CT, this often leads to formulations in which a volume is represented as a collection of slices or patches and a case-level decision is obtained from pooled instance-level information. DR-MIL, for example used a bag-of-slices formulation with deep feature representation and bag-level classification \cite{qi2021drmil}. Xue et al.\ used attention-based multiple-instance learning for CT-based COPD identification and emphasized that disease evidence may be spatially heterogeneous rather than uniformly visible across the volume \cite{xue2023copdmil}. These studies are relevant because they illustrate the methodological design space surrounding the present thesis, coarse labels, unequal instance informativeness, and explicit case-level pooling.

Recent chest-CT work indicates that this remains an active methodological direction. Djahnine et al.\ propose a weakly supervised approach for pathology detection and localization in 3D chest CT, showing that coarse scan-level supervision can still support meaningful abnormality identification and localization in volumetric data \cite{djahnine2024weakly}. Almakky and Yaqub likewise propose a weakly supervised and explainable framework for infection severity classification from volumetric chest CT without relying on dense segmentation labels \cite{almakky2025weakly}. Together, these studies reinforce the broader computational setting of chest CT learning under coarse supervision and heterogeneous volumetric evidence. At the same time, they remain methodologically different from the present thesis, Djahnine et al.\ focus on pathology detection and localization, and Almakky and Yaqub study explainable severity classification, whereas the present work addresses patient-level regression of mucus burden.

The cited studies also differ from the present thesis in an important respect, they address case-level classification or detection-oriented tasks rather than patient-level regression. Their outputs are disease labels, abnormality localization signals, or severity categories rather than continuous burden scores, and their clinical targets differ substantially from mucus plug grading. Their value for the present work is therefore comparative rather than directly transferable. They clarify which weakly supervised aggregation strategies are plausible when a volumetric CT examination must be summarized into one prediction, but they do not directly solve the task considered in this thesis.


\subsection{Transfer Learning in Limited-Data CT Settings}
A second methodological theme relevant to the present thesis concerns transfer learning under limited data. In many medical-imaging settings, the available target cohort is too small to support robust end-to-end feature learning from scratch, which has made pretrained backbones an important practical and methodological tool \cite{kim2022transfer}. This is particularly relevant here because the thesis operates under a small-$N$ constraint and studies whether compact pretrained feature extractors can support patient-level prediction from chest CT under weak supervision.

This literature does not directly solve the present task either. Most transfer-learning studies in medical imaging focus on classification rather than regression, and many operate in larger datasets or with target definitions that differ from patient-level mucus-burden scoring. Nevertheless, the transfer-learning literature is relevant because it provides the conceptual basis for using pretrained backbones and domain-adaptation strategies when task-specific labels are limited. In the present thesis, these ideas are not used as generic performance tricks, but as part of a constrained design space in which model capacity, data availability, and supervision level must be balanced carefully.


\section{Research Gaps and Positioning of This Thesis}

The literature reviewed above suggests that the present thesis sits at the intersection of two partially connected research areas. On one side, CT-based mucus-plug studies have shown that patient-level mucus burden can be scored visually and that this burden carries clinically relevant information \cite{dunican2018mucus,dunican2018autopsy}. More recent work has extended this area in two directions: toward contemporary clinical and pathophysiologic study of mucus plugs \cite{hsieh2025smallairwaymucus}, and toward automated CT-based quantification, validation, detection, and segmentation \cite{vanderveer2025aicopdmucus,raut2025pragmaai,pu2026simulation}. On the other side, weakly supervised and multiple-instance CT learning studies have developed computational strategies for learning from coarse case-level labels and pooled local evidence \cite{qi2021drmil,xue2023copdmil,almakky2025weakly}. What remains less clearly addressed is the combination of these two perspectives: computational prediction of patient-level mucus burden from chest CT under weak supervision and limited data.

This gap is meaningful because the two literatures solve different parts of the problem. The CT mucus-plug literature establishes that the target is clinically interpretable, but it does not address how that target should be estimated automatically from CT. More recent automated mucus-plug studies move closer to computation, but they are still primarily concerned with quantification, validation, detection, or segmentation rather than weakly supervised patient-level regression \cite{vanderveer2025aicopdmucus,raut2025pragmaai,pu2026simulation}. By contrast, the weakly supervised CT literature offers relevant methodological patterns for learning from volumetric image evidence under coarse supervision, but it is mostly focused on case-level classification rather than regression of a clinically defined burden score \cite{qi2021drmil,xue2023copdmil,almakky2025weakly}. Existing work therefore does not directly address the setting studied in this thesis: patient-level mucus score prediction from chest CT when only coarse patient-level labels are available.

A second gap concerns the combination of supervision constraints and data scale. Many related medical-imaging studies either address different targets or operate in settings with annotation structures, data regimes, or model objectives that do not match patient-level mucus burden scoring. The present thesis instead studies a constrained setting in which the target is clinically meaningful, the supervision is weak, and the available cohort is small. This makes the problem distinct from both localized detection or segmentation tasks and more standard case-level classification settings.

The present thesis is therefore positioned as a machine-learning systems study built around an existing clinically meaningful target. It does not introduce a new mucus-plug scoring system, and it does not attempt plug-level detection or segmentation. Instead, it asks how far a slice-based, weakly supervised CT pipeline can go when only patient-level mucus scores are available. This makes the contribution mainly methodological and empirical, the thesis evaluates a constrained prediction pipeline and analyzes which design choices appear useful under limited-data conditions.

\chapter{Theory and Problem Formulation}
\label{chap:theory}

This chapter presents the concepts needed to formulate CT-based mucus-plug prediction as a machine-learning problem. A chest CT examination is treated as volumetric image data composed of multiple local observations, while the target is defined at patient level. This creates a weakly supervised regression setting in which local image evidence must be represented, modelled, and related to one scan-level output.


\section{CT-Based Prediction Problem}
\label{sec:ct_based_prediction_problem}

\subsection{Volumetric CT as Structured Input Data}
\label{sec:volumetric_ct_input}

CT-based prediction problems differ from standard two-dimensional image prediction tasks because the input is a volumetric medical examination rather than a single isolated image. A chest CT scan consists of a sequence of cross-sectional slices that together represent thoracic anatomy \cite{whiting2015ctbasic,dejong2005ctairways}. For a patient or examination indexed by \(i\), the CT volume is represented as

\[
X_i = \{x_{i,1}, x_{i,2}, \dots, x_{i,n_i}\},
\]

where \(X_i\) denotes the full CT examination, \(x_{i,j}\) denotes the \(j\)-th axial slice or local image instance, and \(n_i\) is the number of slices in the volume. Relevant information may be distributed across several slices, localized to a small part of the volume, or absent from large regions of the scan \cite{xue2023copdmil,djahnine2024weakly}. The scan must therefore be treated as a structured set of local observations rather than as one flat image.

\subsection{Patient-Level Targets and Regression Formulation}
\label{sec:patient_level_regression}

In CT-based mucus-plug assessment, the clinically meaningful quantity is commonly summarized as an overall burden at the patient or scan level, rather than as an independent label for each individual slice \cite{dunican2018mucus,dunican2018autopsy}. Let

\[
y_i \in \mathbb{R}
\]

denote the patient-level target associated with examination \(X_i\). The prediction task can then be written as

\[
\hat{y}_i = f(X_i),
\]

where \(f(\cdot)\) is a prediction function that maps the CT examination to an estimated patient-level output \(\hat{y}_i\). Since the target represents a burden measure rather than a discrete class label, the problem is naturally formulated as regression:

\[
f : X_i \mapsto \hat{y}_i \in \mathbb{R}.
\]

This formulation defines both supervision and evaluation at the scan level. Individual slices may contain useful local evidence, but the final prediction is meaningful only as a patient-level quantity.


\subsection{Weak Supervision and Label Granularity}
\label{sec:weak_supervision_label_granularity}
A central theoretical difficulty is the mismatch between the level at which image evidence is observed and the level at which supervision is available. The training data can be written as

\[
\mathcal{D} = \{(X_i, y_i)\}_{i=1}^{N},
\]

where each CT examination \(X_i\) is associated with one patient-level target \(y_i\). However, corresponding slice-level labels

\[
y_{i,j}
\]

are not observed for the individual slices \(x_{i,j}\). The model therefore has access to the overall target for the examination, but not to annotations indicating which slices, regions, or anatomical locations are responsible for that target. This makes the problem weakly supervised, the labels are coarser than the image evidence from which the model must learn \cite{ren2023weakly}. Related CT studies address similar situations by treating a scan as a collection of local instances and combining instance-level evidence into a case-level prediction \cite{qi2021drmil,xue2023copdmil}. The same issue arises here, although the target is a continuous burden score rather than a class label.


\subsection{Implications for CT-Based Learning}
\label{sec:implications_ct_learning}

The task can therefore be understood as patient-level regression under weak supervision. It requires a suitable CT representation, feature extraction from local image evidence, aggregation from slice-level information to a scan-level output, and evaluation at patient level. Limited cohort size makes these choices more important, because highly flexible models may not be supported by the available supervision \cite{dhiman2023samplesize,litjens2017survey}. For volumetric medical imaging, the choice between slice-based, instance-based, and fully volumetric formulations is shaped by annotation level, computational cost, transfer-learning possibilities, and generalization risk \cite{wang2023twoseventyfive}.








\section{Preprocessing of CT Data}
\label{sec:preprocessing_ct_data}

\subsection{Purpose of Preprocessing in CT-Based Learning}
\label{sec:purpose_preprocessing_ct}

Preprocessing transforms CT data into a form suitable for computational modelling. In CT-based learning, this is needed because medical images may vary in scanner settings, acquisition protocols, reconstruction choices, and anatomical coverage \cite{rizzo2018radiomics,litjens2017survey}.

At a general level, preprocessing can be written as a transformation
\[
\tilde{X}_i = T(X_i),
\]
where \(X_i\) denotes the original CT examination, \(T(\cdot)\) denotes a preprocessing operation or sequence of operations, and \(\tilde{X}_i\) denotes the transformed examination used in later analysis. When preprocessing is considered at slice level, the same idea can be written as
\[
\tilde{x}_{i,j} = T(x_{i,j}),
\]
where \(x_{i,j}\) is an axial CT slice and \(\tilde{x}_{i,j}\) is the corresponding preprocessed slice.

The purpose of preprocessing is not to solve the prediction problem directly.
Rather, preprocessing defines how the raw image evidence is made available for later analysis. It may reduce technical variation, impose numerical consistency, or emphasize image structures that are relevant for subsequent representation learning. In digital image processing, such operations are commonly used to improve the suitability of image data for later analysis, for example by adjusting intensity ranges, reducing unwanted variation, or emphasizing selected spatial structures \cite{gonzalez2018digital}.


\subsection{Intensity Transformation and CT Windowing}
\label{sec:intensity_windowing_ct}

CT images are represented in terms of attenuation values, commonly expressed in Hounsfield units \cite{whiting2015ctbasic}. These values cover a wide numerical range, so a restricted intensity mapping is often needed before visual interpretation or computational analysis.

Windowing is a standard way of mapping CT attenuation values into a more focused intensity range. In clinical CT interpretation, different windows are used to emphasize different anatomical or tissue relationships, such as air-filled lung structures or denser soft-tissue regions \cite{whiting2015ctbasic}. More generally, windowing can be understood as an intensity transformation applied before later analysis. For a slice \(x_{i,j}\), a windowed representation can be written as
\[
x_{i,j}^{(w)} = T_w(x_{i,j}),
\]
where \(T_w(\cdot)\) denotes the intensity transformation associated with window setting \(w\).

Windowing does not change the anatomy in the slice; it changes which attenuation relationships are emphasized. This matters in chest CT because lung parenchyma, airways, vessels, and soft-tissue structures are not equally visible under one intensity mapping \cite{dejong2005ctairways}. From a machine-learning perspective, CT windowing is therefore part of the input representation problem. Prior work on multi-window CT learning supports the general idea that different CT windows can provide complementary information, although the choice remains task dependent \cite{lu2020multiwindow}.


\subsection{Spatial Standardization and Slice-Level Inputs}
\label{sec:spatial_standardization_slice_inputs}

CT examinations can also differ in matrix size, field of view, slice thickness, slice spacing, number of slices, and anatomical coverage \cite{rizzo2018radiomics}. These differences create challenges for models that require inputs with consistent numerical dimensions and comparable spatial organization \cite{litjens2017survey}.

Spatial preprocessing addresses this issue by transforming the image data into a more standardized form. Depending on the modelling formulation, this may involve operations such as resizing, resampling, cropping, padding, or restricting the analysis to a defined image region. Abstractly, a spatial preprocessing operation can be included in the same transformation notation
\[
\tilde{x}_{i,j} = T_{\mathrm{spatial}}(x_{i,j}),
\]
where \(T_{\mathrm{spatial}}(\cdot)\) denotes a transformation that changes the spatial format of the slice while aiming to preserve relevant anatomical content.

Different modelling choices require different spatial preparation. A volumetric model, a slice-based model, and a patch-based model may not use the same input units, but each requires image data that can be processed consistently by the learning method \cite{wang2023twoseventyfive}. In slice-level formulations, axial slices act as local observations within the larger examination. This does not make the final task slice-level prediction; supervision remains defined for the full examination \cite{ren2023weakly,xue2023copdmil}.

\subsection{Candidate Slice Relevance}
\label{sec:candidate_slice_relevance}

A volumetric CT examination may contain many slices, but not all slices are equally informative for a scan-level prediction task. This is particularly important in weakly supervised settings, where the label describes the case as a whole but does not identify which local observations are responsible for that label \cite{ren2023weakly}. At a general level, one may describe a candidate set of local observations as
\[
C_i \subseteq X_i,
\]
where \(C_i\) denotes the set of slices or image instances retained as candidate evidence from the full CT examination \(X_i\). This notation is deliberately general: it does not specify how the candidate set is obtained, only that a CT-based learning problem may involve deciding which parts of the volume are considered suitable for later analysis.

Candidate selection can reflect anatomical relevance, image suitability, and redundancy reduction. Related weakly supervised and multiple-instance CT studies highlight that local evidence in volumetric images is often heterogeneous rather than uniformly informative \cite{xue2023copdmil,djahnine2024weakly}. The exact method used to filter or select slices is methodological; the theoretical point is that a volumetric scan contains many local observations whose relevance to the patient-level target may vary.



\section{Feature Extraction from CT Images}
\label{sec:feature_extraction_ct_images}

\subsection{Purpose of Feature Extraction}
\label{sec:purpose_feature_extraction}

Feature extraction transforms image data into representations that make relevant visual information available for prediction. In CT-based learning, the raw image array is high-dimensional, so a feature extractor maps a preprocessed image or region into a more compact feature space.

For a preprocessed CT slice \(\tilde{x}_{i,j}\), feature extraction can be written generally as
\[
z_{i,j} = \phi(\tilde{x}_{i,j}),
\]
where \(\phi(\cdot)\) denotes a feature-extraction function and \(z_{i,j}\) is the resulting feature representation. The notation is deliberately general. The function \(\phi(\cdot)\) may represent a handcrafted descriptor, a learned neural-network mapping, or another transformation that converts image content into a representation suitable for modelling.

Preprocessing prepares the image data, while feature extraction defines which visual properties are made available for prediction. In scan-level CT prediction, useful features should preserve local evidence that may contribute to the patient-level target while reducing variation that is unlikely to help prediction \cite{rizzo2018radiomics,litjens2017survey}.



\subsection{Handcrafted Image Features}
\label{sec:handcrafted_image_features}

Handcrafted image features are predefined descriptors that encode visual properties such as intensity, texture, edges, gradients, local patterns, or spatial relationships. In medical imaging, such features have often been used to describe image structure compactly, especially when labelled data are limited \cite{castellano2004texture,rizzo2018radiomics}.

A general handcrafted feature extractor can be written as
\[
z_{i,j}^{(\mathrm{hand})} = \phi_{\mathrm{hand}}(\tilde{x}_{i,j}),
\]
where \(\phi_{\mathrm{hand}}(\cdot)\) denotes a predefined feature function. The resulting vector \(z_{i,j}^{(\mathrm{hand})}\) may encode local texture, intensity statistics, edge information, or other manually specified image characteristics. Such features are not optimized end-to-end for the target task, but they may provide useful summaries of visual structure when the relevant image patterns are expected to be local or when the available training data are limited.

Their strength is that they encode prior assumptions about useful image properties, which can be valuable with limited labels. Their limitation is that the representation is fixed before learning; if the descriptors miss the relevant evidence, the downstream model cannot recover it.



\subsection{Learned Feature Representations}
\label{sec:learned_feature_representations}

Learned feature representations are estimated from data rather than defined manually. In deep learning, feature extraction becomes part of the learning process rather than a separate fixed descriptor stage \cite{lecun1998gradient,litjens2017survey}.

A learned feature extractor can be written as
\[
z_{i,j}^{(\mathrm{learned})} = \phi_{\theta}(\tilde{x}_{i,j}),
\]
where \(\phi_{\theta}(\cdot)\) is a parameterized mapping with trainable parameters \(\theta\). Unlike handcrafted descriptors, the representation is shaped by the data, the objective function, and the available supervision.

For CT images, learned representations are useful because relevant evidence may be complex, spatially distributed, and difficult to specify manually. Their flexibility is also their main risk: it must be supported by enough data and appropriate supervision, as discussed in Section~\ref{sec:model_capacity_limited_data}.




\subsection{Feature Dimensionality and Representation Quality}
\label{sec:feature_dimensionality_representation_quality}

The quality of a feature representation depends on both its information content and its dimensionality. A useful representation should preserve target-relevant image properties while reducing redundant or unrelated variation. This balance is especially important in medical imaging, where image data are high-dimensional and patient cohorts are often limited \cite{mwangi2014feature,rizzo2018radiomics}.

If a representation is too low-dimensional, relevant image information may be lost. If it is too high-dimensional relative to the number of labelled cases, the downstream model may become more vulnerable to overfitting and unstable estimation. Dimensionality reduction and feature selection can therefore be understood as general strategies for improving the match between the information content of the representation and the amount of supervision available for learning \cite{mwangi2014feature}.

Handcrafted and learned representations encode image information according to different assumptions and may capture complementary aspects of medical images \cite{jeong2024handcrafteddeep,lai2018featurefusion}. Whether this improves prediction is empirical, but conceptually it shows why feature extraction determines what kind of evidence later model stages can use.



\section{CNN Modelling and Transfer Learning}
\label{sec:cnn_modelling_transfer_learning}

This section introduces CNNs, transfer learning, and fine-tuning as concepts needed to understand the implemented CT prediction pipeline under limited supervision.


\subsection{Convolutional Neural Networks for CT Images}
\label{sec:cnn_ct_images}

Convolutional neural networks are neural models designed for structured image data. They apply learnable filters across local image regions, allowing similar visual patterns to be detected across spatial locations \cite{lecun1998gradient,litjens2017survey}.

For a preprocessed CT slice \(\tilde{x}_{i,j}\), a CNN-based mapping can be written abstractly as
\[
z_{i,j} = \phi_{\theta}(\tilde{x}_{i,j}),
\]
where \(\phi_{\theta}(\cdot)\) denotes a convolutional feature extractor with trainable parameters \(\theta\), and \(z_{i,j}\) denotes the resulting learned representation. A task-specific prediction function may then map this representation to an instance-level output,
\[
s_{i,j} = h_{\theta}(z_{i,j}).
\]
This notation separates representation learning from prediction without committing to a specific architecture.

CNNs are relevant for CT images because useful evidence may be local, spatially structured, and difficult to summarize using global intensity statistics alone. However, CNN modelling does not remove the weakly supervised nature of the task: local representations still have to be related to patient-level targets through aggregation and evaluation.


\subsection{Residual Learning}
\label{sec:residual_learning}

Residual learning allows a network block to learn a residual transformation relative to its input rather than a completely new representation from scratch \cite{he2016residual}. This is relevant because deeper CNNs can represent richer features but may also be harder to optimize.

A residual block can be written abstractly as
\[
u_{l+1} = u_l + F_l(u_l; \theta_l),
\]
where \(u_l\) is the input representation at layer or block \(l\), \(F_l(\cdot)\) is a learnable residual function, \(\theta_l\) denotes the parameters of that block, and \(u_{l+1}\) is the resulting output representation. The additive identity connection allows information from earlier layers to pass forward directly, while the residual function learns how the representation should be modified.

For this thesis, residual learning matters mainly as the architectural principle behind the ResNet backbone used later. The choice of a concrete backbone remains a methodological decision.


\subsection{Model Capacity in Limited-Data Settings}
\label{sec:model_capacity_limited_data}

Model capacity refers to the flexibility of a model to represent complex functions. In CNNs, capacity is influenced by network depth, number of parameters, feature dimensionality, and prediction-head complexity. Higher-capacity models can represent more complex relationships, but require sufficient supervision to estimate their parameters reliably.

This issue is particularly important in medical imaging, where labelled datasets are often smaller than in general computer vision and where labels may be expensive or difficult to obtain \cite{litjens2017survey}. At a general level, model learning can be expressed as
\[
\theta^{*} = \arg\min_{\theta} \mathcal{L}_{\mathrm{train}}(\theta).
\]
When the model is highly flexible and the dataset is small, minimizing the training loss alone may not ensure good generalization.

Insufficient sample size can lead to unstable estimates and overly optimistic performance assessment \cite{dhiman2023samplesize}. In volumetric medical imaging, this concern is amplified by high-dimensional input and by the choice between two-dimensional, slice-based, and fully volumetric formulations \cite{wang2023twoseventyfive}. Model capacity is therefore part of the theoretical constraint imposed by limited medical image data.


\subsection{Transfer Learning as Representation Reuse}
\label{sec:transfer_learning_representation_reuse}

Transfer learning addresses limited-data settings by reusing information learned from a source task or source dataset. Instead of initializing all model parameters randomly, a model can begin from parameters already learned elsewhere \cite{kim2022transfer}.

In general notation, pretraining can be described as learning parameters on a source dataset,
\[
\theta_{\mathrm{source}}^{*}
=
\arg\min_{\theta}
\mathcal{L}_{\mathrm{source}}(\theta),
\]
and then using those parameters to initialize learning on the target task,
\[
\theta_0 \leftarrow \theta_{\mathrm{source}}^{*}.
\]
Transfer learning can therefore be understood as representation reuse: the pretrained model provides an initial feature space that can be adapted to the target prediction problem.

Transfer learning is relevant in medical imaging because target datasets are often limited, while deep models may require many labelled examples \cite{litjens2017survey,kim2022transfer}. Pretraining may use natural-image data or medical-image data; RadImageNet is one radiology-domain example \cite{mei2022radimagenet}. Its usefulness depends on the relationship between the source and target domains, so it should be understood as a strategy for limited supervision rather than a guaranteed improvement.


\subsection{Fine-Tuning and Domain Adaptation}
\label{sec:fine_tuning_domain_adaptation}

Fine-tuning is the process of adapting a pretrained model to a target task. After initialization from source-domain parameters, the model is trained further using target-domain data. This can be written abstractly as
\[
\theta_{\mathrm{target}}^{*}
=
\arg\min_{\theta}
\mathcal{L}_{\mathrm{target}}(\theta; \theta_0),
\]
where \(\theta_0\) denotes the pretrained initialization and \(\mathcal{L}_{\mathrm{target}}\) is the target-task objective.

Fine-tuning strategies differ in how much of the pretrained representation is reused. Keeping most layers fixed treats the model mainly as a feature extractor, while updating more layers allows stronger adaptation to the target domain. This balances stability and flexibility: preserving pretrained features may reduce overfitting when target data are limited, while adapting more parameters may be needed when source and target domains differ \cite{kim2022transfer}. The exact fine-tuning procedure is methodological; the theoretical role is to align a reused representation with the target prediction problem.



\section{CT Slice Aggregation}
\label{sec:ct_slice_aggregation}

\subsection{From Local Slice Evidence to Scan-Level Output}
\label{sec:local_slice_to_scan_output}

In scan-level CT prediction, the target is defined for the full examination, while much of the image evidence is observed locally at slice or region level. Aggregation is therefore required to combine local evidence into a single patient-level prediction. Related CT learning work uses similar formulations by treating an examination as a collection of slice-level instances that support a case-level output \cite{qi2021drmil,xue2023copdmil}.

Let \(s_{i,j}\) denote an instance-level score or evidence value derived from slice \(x_{i,j}\). A general aggregation function can be written as
\[
\hat{y}_i = A(s_{i,1}, s_{i,2}, \dots, s_{i,n_i}),
\]
where \(A(\cdot)\) maps local evidence values to a scan-level prediction \(\hat{y}_i\). If only a candidate subset of local observations is considered, the same idea can be written as
\[
\hat{y}_i = A\bigl(\{s_{i,j} : x_{i,j} \in C_i\}\bigr),
\]
where \(C_i \subseteq X_i\) denotes the local observations considered for prediction.

Different aggregation functions express different assumptions about how relevant evidence is distributed across the scan. In volumetric CT, disease-related evidence may be spatially heterogeneous rather than uniformly present throughout the volume \cite{xue2023copdmil}. Aggregation is therefore a conceptual mechanism for relating local CT evidence to a scan-level target.


\subsection{Multiple-Instance Learning Perspective}
\label{sec:mil_perspective}

The aggregation problem can be interpreted through multiple-instance learning. In this view, an input case is a bag of instances, while supervision is available only at bag level. For CT, the scan is the bag and slices or local regions are instances. The available label \(y_i\) describes the bag as a whole, while instance-level labels are unobserved.

Using this terminology, the CT examination
\[
X_i = \{x_{i,1}, x_{i,2}, \dots, x_{i,n_i}\}
\]
is a bag of local observations. Multiple-instance CT studies use this type of formulation to relate local image information to case-level outputs \cite{qi2021drmil,xue2023copdmil}. Although many multiple-instance studies address classification, the same conceptual structure applies to regression: local evidence must be combined into a continuous scan-level estimate rather than a discrete class probability.

This perspective separates two questions: feature extraction concerns how local evidence is represented, while aggregation concerns how that evidence is combined into a scan-level output.


\subsection{Aggregation Assumptions}
\label{sec:aggregation_assumptions}

Aggregation functions encode assumptions about how relevant evidence is distributed across the scan. A diffuse-evidence assumption treats many local observations as informative, while a focal-evidence assumption treats the target as driven mainly by a smaller number of highly informative observations \cite{qi2021drmil,xue2023copdmil}. Mixed or weighted assumptions allow both broad and localized evidence to influence the final prediction. The exact aggregation rule is methodological, but scan-level prediction always requires some assumption about how local CT evidence relates to the patient-level target.


\subsection{Aggregation under Weak Supervision}
\label{sec:aggregation_weak_supervision}

Aggregation becomes especially important under weak supervision because the model does not receive direct feedback about which slices are relevant. An accurate scan-level prediction does not necessarily mean that the model has identified the clinically relevant slices or regions. Different local evidence patterns may lead to similar scan-level outputs, so aggregation supports prediction but does not by itself solve localization or explanation \cite{djahnine2024weakly,almakky2025weakly}. Specific aggregation functions, hyperparameters, and empirical comparisons belong to the methodology and experiments.



\section{Regression and Evaluation}
\label{sec:regression_evaluation}

\subsection{Continuous Target Prediction}
\label{sec:continuous_target_prediction}

The prediction task is formulated as regression because the target is a numerical burden measure rather than a discrete class label. For a CT examination \(X_i\), a parameterized model produces a patient-level prediction
\[
\hat{y}_i = f_{\theta}(X_i),
\]
where \(y_i \in \mathbb{R}\) denotes the corresponding target. The prediction error for patient \(i\) can be written as
\[
e_i = y_i - \hat{y}_i.
\]

This formulation preserves the ordering and magnitude of the target. Errors are interpreted by their numerical size rather than only as correct or incorrect, which is appropriate when burden magnitude is meaningful.


\subsection{Regression Loss Functions}
\label{sec:regression_loss_functions}

Regression models are trained by minimizing a loss function that penalizes disagreement between observed and predicted target values:
\[
\theta^{*}
=
\arg\min_{\theta}
\mathcal{L}(\theta).
\]
A common objective for continuous prediction is mean squared error,
\[
\mathcal{L}_{\mathrm{MSE}}(\theta)
=
\frac{1}{N}
\sum_{i=1}^{N}
\left(y_i - \hat{y}_i\right)^2.
\]

The squared-error form penalizes larger deviations more strongly than smaller deviations. Other regression losses and metrics emphasize different aspects of error, so the optimization loss and the reported evaluation metrics should be distinguished \cite{terven2025lossmetrics,kocak2025medicalmetrics}.


\subsection{Patient-Level Evaluation}
\label{sec:patient_level_evaluation}

Evaluation must be performed at the same level as the target. In scan-level CT prediction, performance is assessed by comparing patient-level predictions with patient-level targets,
\[
\{(y_i, \hat{y}_i)\}_{i=1}^{N}.
\]
Intermediate slice-level values may be useful inside the model, but they should not be treated as independently validated outputs unless corresponding slice-level labels are available.

Error-based metrics summarize the magnitude of prediction errors, while association-based metrics describe how predictions vary with observed targets. Correlation coefficients measure association between continuous variables, but they do not by themselves describe absolute error or calibration \cite{schober2018correlation}. Patient-level regression evaluation should therefore use metrics that reflect the intended meaning of prediction error.


\subsection{Validation and Generalization in Small Cohorts}
\label{sec:validation_generalization_small_cohorts}

Validation estimates how well a model is expected to perform on unseen cases. Performance should be evaluated on data not used to fit the model parameters. If \(\mathcal{D}_{\mathrm{train}}\) denotes the training data and \(\mathcal{D}_{\mathrm{val}}\) denotes held-out validation data, learning and validation can be separated as
\[
\theta^{*}
=
\arg\min_{\theta}
\mathcal{L}_{\mathrm{train}}(\theta; \mathcal{D}_{\mathrm{train}}),
\]
with performance estimated on \(\mathcal{D}_{\mathrm{val}}\).

This separation is especially important in small-cohort medical imaging, where high-dimensional inputs and flexible models can produce overly optimistic performance estimates if evaluation is not carefully separated from training \cite{dhiman2023samplesize,litjens2017survey}. Cross-validation and related resampling strategies can make more efficient use of limited data, but they do not remove the uncertainty caused by small sample size.

Generalization also depends on whether future data resemble the development data. CT examinations may vary across scanners, acquisition protocols, reconstruction settings, and patient populations \cite{rizzo2018radiomics}. Evaluation in a small cohort should therefore be interpreted as evidence about performance under the available data conditions, not as proof of clinical robustness. The theoretical point is that regression evaluation must be patient-level and that generalization claims require caution under limited data.

\chapter{Methodology}
\label{chap:methodology}

This chapter describes how the conceptual formulation from Chapter~\ref{chap:theory} was implemented. It specifies the data representation, slice-selection procedure, model components, aggregation rules, training setup, and reproducibility details used in this thesis.


\section{Overview of the Proposed Pipeline}
\label{sec:overview_of_the_proposed_pipeline}
The implemented pipeline operates at patient level. Each patient \(i\) is represented by a volumetric chest CT examination \(X_i\) and a corresponding target \(y_i\). The volume is preprocessed, reduced to representative slices, converted into model inputs, passed through feature-extraction and prediction modules, and finally aggregated into one patient-level estimate.

Figure~\ref{fig:proposed_model_architecture} gives an overview of the proposed ResNet18-3CH+BoVW pipeline. The individual stages shown in the figure are described in detail in the following sections.

\begin{figure}[H]
\centering

\definecolor{ctBlue}{RGB}{221,235,247}
\definecolor{ctBlueLine}{RGB}{88,126,166}
\definecolor{prepTeal}{RGB}{218,241,238}
\definecolor{prepTealLine}{RGB}{73,139,130}
\definecolor{sliceOrange}{RGB}{254,232,200}
\definecolor{sliceOrangeLine}{RGB}{190,121,51}
\definecolor{reprYellow}{RGB}{255,247,188}
\definecolor{reprYellowLine}{RGB}{169,145,44}
\definecolor{cnnGreen}{RGB}{223,240,216}
\definecolor{cnnGreenLine}{RGB}{85,139,78}
\definecolor{bovwPurple}{RGB}{232,225,245}
\definecolor{bovwPurpleLine}{RGB}{117,91,154}
\definecolor{fusionGray}{RGB}{226,226,226}
\definecolor{fusionGrayLine}{RGB}{95,95,95}
\definecolor{predRed}{RGB}{248,218,218}
\definecolor{predRedLine}{RGB}{166,83,83}

\resizebox{0.92\textwidth}{!}{%
\begin{tikzpicture}[
    font=\scriptsize,
    node distance=5.5mm,
    >={Latex[length=2mm]},
    pipeline/.style={
        draw,
        rounded corners=2pt,
        align=center,
        text width=6.7cm,
        inner sep=4pt,
        minimum height=8mm
    },
    branch/.style={
        draw,
        rounded corners=2pt,
        align=center,
        text width=4.45cm,
        inner sep=4pt,
        minimum height=8mm
    },
    note/.style={
        draw,
        rounded corners=2pt,
        dashed,
        fill=white,
        align=left,
        text width=3.85cm,
        inner sep=4pt
    },
    arrow/.style={->, line width=0.65pt, draw=black!70},
    notearrow/.style={->, dashed, line width=0.5pt, draw=black!55},
    cticon/.pic={
        \draw[fill=white, draw=ctBlueLine, rounded corners=0.6pt, line width=0.35pt]
            (-0.42,-0.30) rectangle (0.42,0.30);
        \draw[fill=black!78, draw=black!55, line width=0.35pt]
            (0,0) circle (0.23);
        \draw[fill=black!32, draw=black!25, line width=0.25pt]
            (-0.05,0.02) ellipse (0.12 and 0.08);
        \draw[fill=black!38, draw=black!25, line width=0.25pt]
            (0.10,-0.04) ellipse (0.10 and 0.07);
        \draw[draw=white, line width=0.35pt]
            (-0.15,0.13) .. controls (-0.02,0.23) and (0.13,0.18) .. (0.17,0.05);
    }
]

\node[pipeline, fill=ctBlue, draw=ctBlueLine, minimum height=16mm, text width=7.0cm] (volume) {
    \hspace*{1.45cm}\parbox{5.2cm}{\centering
    \textbf{Patient CT examination}\\
    Stored volume \(X_i\) and patient-level target \(y_i\)}
};

\node[anchor=center] at ([xshift=-2.75cm]volume.center)
{\includegraphics[height=9mm]{figures/rawct1.png}};



\node[pipeline, below=of volume, fill=prepTeal, draw=prepTealLine] (preprocess) {
    \textbf{Data preprocessing}\\
    HU clipping/canonicalization; slice-wise resizing to \(256 \times 256\)
};

\node[pipeline, below=of preprocess, fill=sliceOrange, draw=sliceOrangeLine] (selection) {
    \textbf{Representative slice selection}\\
    Central restriction; informativeness filtering; PCA descriptors;\\
    MiniBatchKMeans diversity sampling; selected set \(\mathcal{S}_i\)
};

\node[pipeline, below=of selection, fill=reprYellow, draw=reprYellowLine] (representation) {
    \textbf{Three-channel high-pass representation}\\
    For each selected slice \(j \in \mathcal{S}_i\):\\
    lung-window, mediastinal-window, and high-pass channels\\
    \(r_{i,j}^{(3,\mathrm{HP})}\)

};

\node[note, right=7mm of representation] (hpnote) {
    \textbf{High-pass channel}\\
    Mediastinal source, Gaussian blur subtraction, robust normalization, and lung-mask suppression.
};

\node[branch, below=10mm of representation, xshift=-2.75cm, fill=cnnGreen, draw=cnnGreenLine] (cnn) {
    \textbf{CNN branch}\\
    ResNet18 backbone\\
    learned feature \(z_{i,j}\)
};

\node[branch, below=10mm of representation, xshift=2.75cm, fill=bovwPurple, draw=bovwPurpleLine] (bovw) {
    \textbf{ORB-BoVW branch}\\
    ORB descriptors; BoVW histogram \(b_{i,j}\);\\
    projection to \(\bar{b}_{i,j}\)
};

\node[pipeline, below=35mm of representation, fill=fusionGray, draw=fusionGrayLine] (fuse) {
    \textbf{Feature fusion}\\
    \textit{Proposed ResNet18-3CH+BoVW model}\\
    \(v_{i,j} = [z_{i,j};\bar{b}_{i,j}]\)
};

\node[pipeline, below=of fuse, fill=fusionGray, draw=fusionGrayLine] (head) {
    \textbf{Regression head}\\
    Slice-level score \(s_{i,j}=g_{\omega}(v_{i,j})\)
};

\node[pipeline, below=of head, fill=predRed, draw=predRedLine] (aggregation) {
    \textbf{Patient-level aggregation}\\
    Hybrid aggregation of mean and top-\(k\) slice scores\\
    \(\hat{y}_i=A(\{s_{i,j}:j\in\mathcal{S}_i\})\)
};

\node[pipeline, below=of aggregation, fill=predRed, draw=predRedLine] (prediction) {
    \textbf{Patient-level prediction}\\
    Predicted mucus-burden score \(\hat{y}_i\), compared with \(y_i\)
};

\node[note, right=7mm of aggregation] (weaknote) {
    \textbf{Weak supervision}\\
    Selected slices inherit \(y_i\) during training. Slice scores are intermediate outputs, not verified slice-level mucus labels.
};

\draw[arrow] (volume) -- (preprocess);
\draw[arrow] (preprocess) -- (selection);
\draw[arrow] (selection) -- (representation);
\draw[arrow] (representation.south) -- ++(0,-0.25) -| (cnn.north);
\draw[arrow] (representation.south) -- ++(0,-0.25) -| (bovw.north);
\draw[arrow] (cnn.south) -- ++(0,-0.45) -| (fuse.north);
\draw[arrow] (bovw.south) -- ++(0,-0.45) -| (fuse.north);
\draw[arrow] (fuse) -- (head);
\draw[arrow] (head) -- (aggregation);
\draw[arrow] (aggregation) -- (prediction);
\draw[notearrow] (representation.east) -- (hpnote.west);
\draw[notearrow] (aggregation.east) -- (weaknote.west);

\end{tikzpicture}%
}
\caption[Proposed ResNet18-3CH+BoVW architecture]{Architecture of the proposed ResNet18-3CH+BoVW model used in the internal experiments. Selected CT slices are converted to the three-channel high-pass representation, ResNet18 CNN features and ORB-BoVW handcrafted features are fused at slice level, and the resulting slice-level scores are aggregated into one patient-level mucus-burden prediction. Training is weakly supervised because selected slices inherit the patient-level label, the slice scores are intermediate outputs rather than verified slice-level mucus labels.}
\label{fig:proposed_model_architecture}
\end{figure}



\subsection{Notation}
\label{sec:notation}

This chapter uses the same basic notation as Chapter~3. For patient \(i\), the CT volume is written as
\[
X_i = \{x_{i,1}, x_{i,2}, \dots, x_{i,n_i}\},
\]
where \(x_{i,j}\) is axial slice \(j\), and \(n_i\) is the number of slices in the volume. The patient-level target is written as
\[
y_i \in \mathbb{R}.
\]
Table~\ref{tab:methodology_notation} summarizes the notation used in the methodology. Index sets are distinguished from the image sets they refer to, for example \(\mathcal{S}_i\) contains selected slice indices, while \(S_i\) contains the corresponding selected slices.

For consistency with Chapter~3, all mathematical slice indices are written as \(j \in \{1,\dots,n_i\}\). In the Python implementation, the same slice positions are stored as zero-based array indices.

\begin{table}[H]
\centering
\caption{Notation used in the methodology chapter.}
\label{tab:methodology_notation}
\begin{tabular}{ll}
\toprule
Symbol & Meaning \\
\midrule
\(i\) & Patient or case index \\
\(X_i\) & Full CT volume for patient \(i\) \\
\(x_{i,j}\) & Axial slice \(j\) from volume \(X_i\) \\
\(n_i\) & Number of slices in volume \(X_i\) \\
\(\tilde{x}_{i,j}\) & Preprocessed version of slice \(x_{i,j}\) \\
\(\tilde{X}_i\) & Preprocessed CT volume for patient \(i\) \\
\(y_i\) & Patient-level target \\
\(\hat{y}_i\) & Patient-level prediction \\
\(\mathcal{C}_i\) & Candidate slice index set \\
\(C_i\) & Candidate slice set \\
\(\mathcal{F}_i\) & Filtered slice index set \\
\(F_i\) & Filtered slice set \\
\(\mathcal{G}_i\) & Reserved high-informativeness slice index set \\
\(\mathcal{R}_i\) & Remaining filtered candidate index set for diversity sampling \\
\(\mathcal{D}_i\) & Diversity-selected slice index set \\
\(\mathcal{S}_i\) & Final selected slice index set \\
\(S_i\) & Final selected slice set \\
\(d_{i,j}\) & Low-resolution slice-selection descriptor for slice \(j\) \\
\(p_{i,j}\) & PCA-projected slice-selection descriptor for slice \(j\) \\
\(K_{\mathrm{sel}}\) & Requested number of selected slices per patient \\
\(K_{\mathrm{inf},i}\) & Effective number of reserved informative slices for patient \(i\) \\
\(K_{\mathrm{div},i}\) & Number of diversity-selected slices requested for patient \(i\) \\
\(r_{i,j}\) & Model input representation for slice \(j\) \\
\(z_{i,j}\) & CNN feature vector for slice \(j\) \\
\(b_{i,j}\) & Handcrafted feature vector for slice \(j\), when used \\
\(\bar{b}_{i,j}\) & Projected handcrafted feature vector, when used \\
\(v_{i,j}\) & Fused feature vector for slice \(j\), when used \\
\(s_{i,j}\) & Slice-level prediction score \\
\(P_{\mathrm{prep}}\) & Fixed preprocessing operation \\
\(R\) & Input-representation function \\
\(\Phi_{\psi}\) & CNN feature extractor with parameters \(\psi\) \\
\(g_{\omega}\) & Prediction head with parameters \(\omega\) \\
\(\psi\) & Parameters of the CNN feature extractor \\
\(\omega\) & Parameters of the prediction head \\
\(\eta\) & Parameters of the handcrafted feature projection branch, when used \\
\(\Theta\) & Full set of trainable model parameters \\
\(A\) & Patient-level aggregation rule \\
\(P_{\mathrm{PCA}}\) & PCA projection matrix used during slice selection \\
\(\mu_k\) & Centroid of cluster \(k\) during representative selection \\
\(k\) & Number of top slice scores used in top-\(k\) aggregation \\
\(\lambda\) & Mixing weight used in hybrid aggregation, when applicable \\
\bottomrule
\end{tabular}
\end{table}

The notation keeps fixed operations and trainable model components separate. For example, \(P_{\mathrm{prep}}\) is fixed preprocessing, \(R\) constructs the model input representation, \(\Phi_{\psi}\) extracts CNN features, and \(g_{\omega}\) produces slice-level scores.


\subsection{End-to-End Processing Pipeline}
\label{sec:end_to_end_prediction_chain}

The full pipeline can be summarized as the path from an indexed CT volume to one patient-level prediction:
\[
X_i
\rightarrow
\tilde{X}_i
\rightarrow
\mathcal{S}_i
\rightarrow
\{r_{i,j}: j \in \mathcal{S}_i\}
\rightarrow
\{z_{i,j}: j \in \mathcal{S}_i\}
\rightarrow
\{s_{i,j}: j \in \mathcal{S}_i\}
\rightarrow
\hat{y}_i.
\]
The chain shows the main stages: preprocessing, slice selection, input representation, feature extraction, slice-level prediction, and patient-level aggregation.

The first step is preprocessing. Each raw slice \(x_{i,j}\) is transformed into a preprocessed slice
\[
\tilde{x}_{i,j} = P_{\mathrm{prep}}(x_{i,j}),
\]
where \(P_{\mathrm{prep}}(\cdot)\) denotes the fixed preprocessing operations applied before slice selection and model input construction. The preprocessed volume is then written as
\[
\tilde{X}_i
=
\{\tilde{x}_{i,1}, \tilde{x}_{i,2}, \dots, \tilde{x}_{i,n_i}\}.
\]

Slice selection defines the selected index set
\[
\mathcal{S}_i \subseteq \{1,\dots,n_i\},
\]
and the corresponding selected slice set is
\[
S_i
=
\{\tilde{x}_{i,j} : j \in \mathcal{S}_i\}.
\]
The construction of \(\mathcal{S}_i\) is described in Section~\ref{sec:slice_selection_strategy}. Only these selected slices are passed to the later model stages.

Each selected slice is then converted into a model input representation:
\[
r_{i,j} = R(\tilde{x}_{i,j}),
\qquad j \in \mathcal{S}_i,
\]
where \(R(\cdot)\) denotes the input-representation function. Depending on the experiment, this function produces either a two-window representation or a three-channel representation. The specific representations are described in Section~\ref{sec:input_representation}.

The represented slice is passed through the CNN feature extractor:
\[
z_{i,j} = \Phi_{\psi}(r_{i,j}),
\qquad j \in \mathcal{S}_i,
\]
where \(\Phi_{\psi}(\cdot)\) denotes the CNN backbone with trainable parameters \(\psi\). The prediction head maps this feature vector to a slice-level score,
\[
s_{i,j} = g_{\omega}(z_{i,j}),
\qquad j \in \mathcal{S}_i,
\]
where \(g_{\omega}(\cdot)\) has trainable parameters \(\omega\). In experiments with handcrafted feature fusion, the CNN feature vector is first combined with a projected handcrafted feature vector before prediction. This gives a fused vector
\[
v_{i,j} = [z_{i,j}; \bar{b}_{i,j}],
\]
and the prediction head is applied to \(v_{i,j}\) instead of \(z_{i,j}\).

The slice-level scores are then aggregated into one patient-level prediction:
\[
\hat{y}_i =
A\bigl(\{s_{i,j} : j \in \mathcal{S}_i\}\bigr),
\]
where \(A(\cdot)\) denotes the patient-level aggregation rule. The implemented aggregation variants are described in Section~\ref{sec:patient_level_prediction_aggregation}.

The prediction problem is therefore defined at patient level: the final model output \(\hat{y}_i\) is compared with the selected patient-level target \(y_i\). Conceptually, the desired patient-level regression objective can be written as
\[
\mathcal{L}_{\mathrm{patient}}
=
\frac{1}{|\mathcal{I}_{\mathrm{train}}|}
\sum_{i \in \mathcal{I}_{\mathrm{train}}}
\left(y_i - \hat{y}_i\right)^2.
\]
This expression describes the level at which the task and validation are defined. The implemented optimization uses weak slice-level training: sampled slice scores \(s_{i,j}\) are trained against the corresponding patient label, and patient-level aggregation is applied during validation and model selection. The exact training loss is described in Section~\ref{sec:learning_objective}. Aggregation settings such as \(k\) and \(\lambda\) are fixed experimental settings rather than trainable parameters.



\section{Dataset and Target Definition}
\label{sec:dataset_target_definition}

This section defines the dataset objects used throughout the methodology. Each included case corresponds to one indexed chest CT examination and one patient-level mucus-burden target.


\subsection{Dataset Description}
\label{sec:dataset_description}

The implemented pipeline is organized around a dataset index where each row corresponds to one patient-level case. The index links the patient identifier, stored imaging volume, target variables, and fold information. The original imaging studies come from chest CT examinations stored in DICOM format \cite{dicomstandard2026}, but the model does not read the raw DICOM series directly during training. Instead, each indexed case points to a stored volume through \texttt{volume\_path}.

For patient \(i\), the CT volume is denoted as
\[
X_i = \{x_{i,1}, x_{i,2}, \dots, x_{i,n_i}\},
\]
where \(x_{i,j}\) is the \(j\)-th axial slice and \(n_i\) is the number of slices in the volume. In array form, the stored volume can be written as
\[
X_i \in \mathbb{R}^{n_i \times H_i \times W_i},
\]
where \(H_i\) and \(W_i\) are the in-plane image dimensions before later resizing or model-specific input transformation.


\subsection{Cohort Selection and Inclusion Criteria}
\label{sec:cohort_selection_inclusion_criteria}

The study cohort contains
\[
N = 32
\]
indexed patient cases. Because the cohort is small, unsuitable cases can have a large influence on model fitting, validation, and interpretation.

The final cohort was defined by retaining cases that had both a usable chest CT volume and an associated patient-level mucus-burden target. Let \(\mathcal{I}\) denote the set of included patient indices:
\[
\mathcal{I}
=
\{i : X_i \text{ is a valid chest CT examination and } y_i \text{ is available}\}.
\]
The number of included cases is therefore
\[
|\mathcal{I}| = N.
\]

Preprocessing, slice selection, training, and evaluation are applied only to cases in \(\mathcal{I}\), ensuring that every modelled case has both a valid imaging input and a patient-level target.


\subsection{CT Examination Selection}
\label{sec:ct_examination_selection}

During dataset preparation, the indexed imaging files were checked to ensure that they corresponded to usable chest CT examinations. Records with non-lung content or otherwise unsuitable acquisitions were excluded before preprocessing and model development.

For each retained patient \(i \in \mathcal{I}\), the pipeline uses the stored volume referenced by the corresponding \texttt{volume\_path} entry in the dataset index. This volume defines the imaging input \(X_i\) used in the later stages:
\[
\texttt{volume\_path}_i \longmapsto X_i,
\qquad i \in \mathcal{I}.
\]
This case-level selection is separate from the slice-selection strategy in Section~\ref{sec:slice_selection_strategy}, which is applied later within each included volume.


\subsection{Patient-Level Target Definition}
\label{sec:patient_level_target_definition}

The target variable represents the patient-level mucus-burden score for each included CT examination. The raw target is based on two independent reader assessments:
\[
y_i^{(1)}
\quad \text{and} \quad
y_i^{(2)}
\]
The raw patient-level target is defined as the median of these two scores:
\[
y_i^{(\mathrm{raw})}
=
\operatorname{median}\!\left(y_i^{(1)}, y_i^{(2)}\right).
\]
With two reader scores, this median corresponds to the midpoint when the two assessments differ. This is why the stored raw target can contain intermediate values, even if the individual reader scores are assigned on a discrete scoring scale.

In the dataset index, the resulting raw patient-level target is stored as \texttt{label}. In the included cohort, the observed raw mucus-burden labels satisfy
\[
0.0 \leq y_i^{(\mathrm{raw})} \leq 7.5
\]
The raw target is treated as continuous because it represents burden magnitude rather than class membership.

The dataset index also contains a transformed version of the same target, stored as \texttt{label\_log10}. This stored field is denoted by
\[
y_i^{(\log)}.
\]
The implementation uses the precomputed \texttt{label\_log10} value from the dataset index rather than recomputing the transformation during training. The raw target was used in earlier baseline experiments, while the main internal configuration uses the log-transformed target. When a model is trained on \texttt{label\_log10}, predictions can be transformed back to the raw mucus-score scale before reporting MAE and RMSE.

To keep the later methodology notation compact, \(y_i\) denotes the selected target variable for the experiment being described:
\[
y_i
=
\begin{cases}
y_i^{(\mathrm{raw})}, & \text{raw-label experiments},\\
y_i^{(\log)}, & \text{main log-target configuration}.
\end{cases}
\]
Each included case can therefore be represented as
\[
\left(X_i, y_i^{(\mathrm{raw})}, y_i^{(\log)}, \mathrm{id}_i\right),
\qquad i \in \mathcal{I},
\]
where \(X_i\) is the CT volume, \(y_i^{(\mathrm{raw})}\) is the raw patient-level mucus-burden target, \(y_i^{(\log)}\) is the stored log-transformed target, and \(\mathrm{id}_i\) is the patient or case identifier.



\subsection{External Transfer Dataset: MosMed}
\label{sec:external_transfer_dataset_mosmed}

In addition to the internal mucus-score dataset, MosMedData is used for the transfer-learning experiments. MosMedData is a public chest CT dataset containing 1110 CT scans collected during the COVID-19 epidemic \cite{morozov2020mosmeddata}. In this thesis, it is used as an external transfer dataset, not as external validation for mucus-burden prediction, because its target describes a different case-level CT task.

For a MosMed case \(i\), the CT volume is written as
\[
X_i^{(\mathrm{M})},
\]
with a corresponding case-level target
\[
y_i^{(\mathrm{M})}.
\]
The superscript \((\mathrm{M})\) is used only to indicate that the case belongs to the MosMed dataset.

The MosMed target is an ordinal case-level CT severity label related to COVID-19 lung involvement. In MosMedData, the CT severity categories range from cases without CT signs of viral pneumonia to cases with severe lung involvement \cite{morozov2020mosmeddata}. In the implementation, this case-level label is represented as an ordinal numerical target,
\[
y_i^{(\mathrm{M})} \in \{0,1,2,3,4\},
\]
where lower values correspond to no or limited CT involvement and higher values correspond to more extensive CT involvement.

This target is not a mucus-burden score. MosMed is therefore used only to study representation transfer and domain shift between related CT tasks. The model is trained or evaluated on the MosMed severity target when MosMed is the target domain, and on the internal mucus-burden target when the mucus dataset is the target domain.




\section{Data Preprocessing}
\label{sec:data_preprocessing}

Data preprocessing converts indexed stored CT volumes into HU-like slice arrays used by slice selection and model input construction. For each patient \(i\), the input volume is written as
\[
X_i = \{x_{i,1}, x_{i,2}, \dots, x_{i,n_i}\},
\qquad
X_i \in \mathbb{R}^{n_i \times H_i \times W_i},
\]
where \(n_i\) is the number of axial slices and \(H_i\) and \(W_i\) are the original in-plane dimensions. The preprocessing operation is denoted by \(P_{\mathrm{prep}}\). Applied slice-wise, it gives
\[
\tilde{x}_{i,j} = P_{\mathrm{prep}}(x_{i,j}),
\]
and the corresponding preprocessed volume is
\[
\tilde{X}_i =
\{\tilde{x}_{i,1}, \tilde{x}_{i,2}, \dots, \tilde{x}_{i,n_i}\}.
\]
The operation \(P_{\mathrm{prep}}\) is fixed and contains no trainable model parameters. Window-specific channel construction and resizing to the final CNN input size are described separately in Section~\ref{sec:input_representation}.


\subsection{Volume Loading}
\label{sec:volume_loading}

For each included case \(i\), the workflow loads the CT volume from \texttt{volume\_path}. Let \(p_i\) denote this stored path. The original imaging studies are DICOM examinations \cite{dicomstandard2026}, but the training pipeline does not read the raw DICOM files directly:
\[
X_i = \mathrm{Load}(p_i).
\]
In the internal CT pipeline, these stored volumes are loaded as NumPy arrays using \texttt{np.load}, with memory mapping. The loaded array is interpreted as an ordered axial stack with shape
\[
X_i \in \mathbb{R}^{n_i \times H_i \times W_i}.
\]
The slice \(x_{i,j}\) therefore represents slice \(j\) along the first array dimension. Metadata such as slice count and voxel spacing are stored in the dataset index, but model-input construction operates on the loaded volume array.


\subsection{Intensity Normalization}
\label{sec:intensity_normalization}

The stored internal volumes are treated as HU-like arrays, meaning that the numerical values are interpreted as CT attenuation-style intensities following the Hounsfield scale convention \cite{hounsfield1973computerized}. Before slice-level quality-control metrics and representation-specific helper functions are applied, intensities are clipped to a canonical range:
\[
a_{\mathrm{HU}} = -1350,
\qquad
b_{\mathrm{HU}} = 600
\]
represent the lower and upper canonical clipping bounds. For a slice \(x_{i,j}\), clipping is written as
\[
x^{(\mathrm{clip})}_{i,j}
=
\min\!\left(\max\!\left(x_{i,j}, a_{\mathrm{HU}}\right), b_{\mathrm{HU}}\right).
\]
After clipping, the values are scaled to the interval \([0,1]\) when a normalized helper image is needed:
\[
x^{(\mathrm{norm})}_{i,j}
=
\frac{x^{(\mathrm{clip})}_{i,j} - a_{\mathrm{HU}}}
{b_{\mathrm{HU}} - a_{\mathrm{HU}}}.
\]
This gives slice-selection descriptors and later channel-construction helpers a consistent numerical range. In this chapter, \(\tilde{x}_{i,j}\) denotes the HU-like preprocessed slice after loading and canonical intensity handling. Temporary normalized helper images such as \(x^{(\mathrm{norm})}_{i,j}\) are used inside later selection and representation functions, but they are not the final CNN input representation.


\subsection{Spatial Handling and Input Resizing}
\label{sec:spatial_handling_input_resizing}

The implementation does not perform physical voxel-spacing resampling before slice selection. The axial slice count \(n_i\) is not resampled, and no spacing normalization is applied along the \(z\)-axis before representative slice indices are selected.

The stored HU-like slices are also not resized globally before slice selection. The sampling procedure operates on the loaded volume array. A separate low-resolution descriptor image is created only for PCA and clustering, as described in Section~\ref{sec:feature_extraction_slice_selection}.

The fixed CNN input size is
\[
H_0 = W_0 = 256.
\]
This resizing is applied slice-wise during input-representation construction. After a selected slice index \(j \in \mathcal{S}_i\) has been retrieved, the slice is converted into the model input tensor
\[
r_{i,j}
=
R(\tilde{x}_{i,j}),
\]
where \(R(\cdot)\) includes HU clipping, window normalization, resizing to \(256 \times 256\), channel stacking, and the optional high-frequency channel described in Section~\ref{sec:input_representation}.



\section{Slice Selection Strategy}
\label{sec:slice_selection_strategy}

The CNN is not given every axial slice from a CT volume. Instead, the pipeline selects a representative subset from each preprocessed volume before input-representation construction.

For patient \(i\), the preprocessed CT volume is written as
\[
\tilde{X}_i =
\{\tilde{x}_{i,1}, \tilde{x}_{i,2}, \dots, \tilde{x}_{i,n_i}\},
\]
where the slices are ordered along the axial direction. The selected slice indices are denoted by
\[
\mathcal{S}_i \subseteq \{1,\dots,n_i\}.
\]
The corresponding selected slice set is
\[
S_i =
\{\tilde{x}_{i,j}: j \in \mathcal{S}_i\}.
\]
These selected slices are passed to the input-representation stage described in Section~\ref{sec:input_representation}.



\subsection{Overview of the Selection Pipeline}
\label{sec:overview_selection_pipeline}

Slice selection is treated as a fixed preprocessing step applied separately to each CT volume:
\[
\tilde{X}_i
\rightarrow
\mathcal{C}_i
\rightarrow
\mathcal{F}_i
\rightarrow
\{d_{i,j}\}_{j \in \mathcal{F}_i}
\rightarrow
\{p_{i,j}\}_{j \in \mathcal{F}_i}
\rightarrow
\mathcal{S}_i.
\]
The pipeline first keeps a central candidate set \(\mathcal{C}_i\). It then filters this set using simple quality-control and informativeness checks, giving the filtered index set \(\mathcal{F}_i\). For the remaining slices, lightweight descriptors \(d_{i,j}\) are computed only for the purpose of slice selection. These descriptors are reduced with PCA to obtain \(p_{i,j}\), and MiniBatchKMeans is then used to choose slices that cover different parts of the descriptor space. A small set of high-informativeness slices is reserved before the diversity-based fill step.

This procedure is not trained to detect mucus plugs. It uses no patient targets, slice-level labels, or CNN features; it only chooses a stable and varied set of axial slices before prediction.


\subsection{Central Slice Restriction}
\label{sec:central_slice_restriction}

The first step restricts sampling to the central portion of the axial stack. In the final implementation, the center fraction is
\[
\rho_{\mathrm{center}} = 0.75.
\]
For a volume with \(n_i\) slices, the implemented central restriction corresponds to a symmetric integer margin
\[
m_i =
\left\lfloor
\frac{(1-\rho_{\mathrm{center}})n_i}{2}
\right\rfloor,
\]
with one-based lower and upper limits
\[
j_{\min,i} = m_i + 1,
\qquad
j_{\max,i} = n_i - m_i.
\]
The central candidate index set is therefore
\[
\mathcal{C}_i =
\{j : j_{\min,i} \leq j \leq j_{\max,i}\}.
\]
If the central candidate set contains no valid index, the implementation falls back to the middle slice. If the number of central candidates is already less than or equal to the requested number of representative slices, all central candidates are returned without further reduction.


\subsection{Informativeness Filtering}
\label{sec:informativeness_filtering}

After central restriction, the implementation computes simple per-slice quality-control and informativeness measures from the HU-like slice values. Before these measures are computed, each slice is clipped to \([-1350,600]\). The threshold choices are HU-informed implementation heuristics, not diagnostic thresholds \cite{whiting2015ctbasic,dejong2005ctairways}.

For each candidate slice \(j \in \mathcal{C}_i\), the code computes
\[
\mathrm{lung\_frac}_{i,j},\quad
\mathrm{air\_frac}_{i,j},\quad
\mathrm{hi\_dens\_frac}_{i,j},\quad
\mathrm{body\_frac}_{i,j},\quad
\mathrm{mean\_hu}_{i,j},\quad
\mathrm{std\_hu}_{i,j}.
\]
These quantities are implementation-level descriptors rather than clinical measurements. \(\mathrm{lung\_frac}_{i,j}\) is the fraction below \(-350\) HU and acts as a rough lung-like content proxy. \(\mathrm{air\_frac}_{i,j}\) is the fraction below \(-950\) HU and detects slices that are almost entirely air-like. \(\mathrm{hi\_dens\_frac}_{i,j}\) counts pixels within the rough lung mask above \(-200\) HU, while \(\mathrm{body\_frac}_{i,j}\) is based on the fraction above \(-300\) HU. These cutoffs are used only to remove clearly unhelpful slices and rank filtered slices for reserved inclusion.

The filtering decision is written as a binary function
\[
q(\tilde{x}_{i,j}) \in \{0,1\}.
\]
In the implemented configuration, a slice is retained when
\[
q(\tilde{x}_{i,j}) = 1
\quad \Longleftrightarrow \quad
\mathrm{lung\_frac}_{i,j} \geq 0.06,\quad
\mathrm{air\_frac}_{i,j} \leq 0.98,\quad
\mathrm{body\_frac}_{i,j} \leq 0.45.
\]
The filtered index set is therefore
\[
\mathcal{F}_i =
\{j \in \mathcal{C}_i : q(\tilde{x}_{i,j}) = 1\},
\]
with
\[
F_i =
\{\tilde{x}_{i,j}: j \in \mathcal{F}_i\}.
\]
If no slices pass the filter, the implementation falls back to evenly spaced slices from the central candidate set. If too few slices remain for the requested selection size, all filtered slices are returned without diversity sampling.

The same metrics reserve a subset of high-informativeness slices before diversity sampling. Filtered slices are ranked by \(\mathrm{hi\_dens\_frac}_{i,j}\), and the highest-ranked slices are included before clustering-based fill. The effective number of reserved informative slices is
\[
K_{\mathrm{inf},i}
=
\min\!\left(
N_{\mathrm{informative}},
|\mathcal{F}_i|,
\max(8,\lfloor K_{\mathrm{sel}}/6 \rfloor)
\right).
\]
With \(K_{\mathrm{sel}}=160\), this caps the effective informative subset at \(26\) slices when enough filtered candidates are available, even though \(N_{\mathrm{informative}}=30\).


\subsection{Feature Extraction for Slice Selection}
\label{sec:feature_extraction_slice_selection}

The descriptors used for diversity sampling are separate from CNN features; \(z_{i,j}\) is reserved for CNN feature vectors later in the methodology.

For each filtered slice \(j \in \mathcal{F}_i\), the implementation constructs a lightweight descriptor
\[
d_{i,j}
=
D_{\mathrm{sel}}(\tilde{x}_{i,j}).
\]
The descriptor is formed from a low-resolution version of the HU-like slice. In the implementation, the retained slices are clipped to \([-1350,400]\), mapped to \([0,1]\), downsampled to \(64 \times 64\), and flattened. Thus,
\[
d_{i,j}
=
\operatorname{vec}\!\left(
\operatorname{Down}_{64 \times 64}(\tilde{x}_{i,j})
\right),
\qquad
d_{i,j} \in \mathbb{R}^{4096}.
\]
The descriptor matrix for patient \(i\) is standardized feature-wise before PCA:
\[
\bar{d}_{i,j}
=
\frac{d_{i,j} - \mu_d}{\sigma_d + \epsilon},
\]
where \(\mu_d\) and \(\sigma_d\) are computed over the filtered slices from the same patient volume. These descriptors support only unsupervised representative selection and are not passed to the final CNN regressor.


\subsection{Dimensionality Reduction (PCA)}
\label{sec:dimensionality_reduction_pca}

After standardization, PCA is fitted to the slice-selection descriptors for the current patient volume. Using \(P_{\mathrm{PCA}}\) to denote the fitted projection matrix, the reduced descriptor is written as
\[
p_{i,j}
=
P_{\mathrm{PCA}}^{\top}\bar{d}_{i,j},
\qquad j \in \mathcal{F}_i.
\]
The configured PCA component cap is
\[
N_{\mathrm{PCA}} = 24.
\]
The effective number of components is computed as
\[
r_i =
\min\!\left(
N_{\mathrm{PCA}},
4096,
|\mathcal{F}_i|-1
\right),
\]
so the PCA dimensionality is reduced automatically if too few filtered slices are available. PCA reduces the feature dimension before clustering, consistent with its general role in high-dimensional medical-image feature workflows \cite{mwangi2014feature}. If PCA fitting fails, the implementation falls back to the standardized descriptors \(\bar{d}_{i,j}\) without projection.


\subsection{Diversity Sampling with MiniBatchKMeans}
\label{sec:diversity_sampling_minibatchkmeans}

After filtering and descriptor projection, the remaining step is to complete the selected slice set with slices that cover variation within the filtered candidate pool. Let \(\mathcal{G}_i \subseteq \mathcal{F}_i\) denote the reserved high-informativeness slice indices described in Section~\ref{sec:informativeness_filtering}. The remaining candidates for diversity sampling are
\[
\mathcal{R}_i
=
\mathcal{F}_i \setminus \mathcal{G}_i.
\]
The number of additional slices requested from the diversity-sampling stage is
\[
K_{\mathrm{div},i}
=
K_{\mathrm{sel}} - |\mathcal{G}_i|.
\]
The effective number of diversity clusters is
\[
K^{(\mathrm{eff})}_{\mathrm{div},i}
=
\min\!\left(K_{\mathrm{div},i}, |\mathcal{R}_i|\right).
\]

When \(K^{(\mathrm{eff})}_{\mathrm{div},i} > 0\), MiniBatchKMeans is fitted to the PCA-projected descriptors \(\{p_{i,j}: j \in \mathcal{R}_i\}\). MiniBatchKMeans is a computationally efficient variant of k-means clustering \cite{sculley2010minibatch}. The clustering objective can be written as
\[
\min_{\{\mu_k\}_{k=1}^{K^{(\mathrm{eff})}_{\mathrm{div},i}}}
\sum_{j \in \mathcal{R}_i}
\min_{1 \leq k \leq K^{(\mathrm{eff})}_{\mathrm{div},i}}
\left\|p_{i,j} - \mu_k\right\|_2^2,
\]
where \(\mu_k\) denotes the centroid of cluster \(k\). This step chooses candidates that cover different regions of the projected descriptor space; it does not assign clinical labels to slices.

For each centroid, the representative diversity slice is chosen as the remaining filtered slice whose projected descriptor is closest to that centroid:
\[
j_k^{*}
=
\arg\min_{j \in \mathcal{R}_i}
\left\|p_{i,j} - \mu_k\right\|_2^2,
\qquad
k = 1,\dots,K^{(\mathrm{eff})}_{\mathrm{div},i}.
\]
The resulting diversity-selected indices are
\[
\mathcal{D}_i
=
\{j_1^{*}, j_2^{*}, \dots, j_{K^{(\mathrm{eff})}_{\mathrm{div},i}}^{*}\}.
\]
The same slice can be nearest to more than one centroid, so \(\mathcal{D}_i\) may contain duplicate indices before post-processing.

The final selected index set is constructed by combining the reserved informative indices with the diversity-selected indices:
\[
\mathcal{S}_i
=
\operatorname{EnsureUnique}_{K_{\mathrm{sel}}}
\left(
[\mathcal{G}_i,\mathcal{D}_i]
\right).
\]
Here, \([\mathcal{G}_i,\mathcal{D}_i]\) is the ordered candidate list, with reserved informative slices placed before diversity-selected slices. The \(\operatorname{EnsureUnique}_{K_{\mathrm{sel}}}(\cdot)\) helper removes duplicate indices, keeps only valid filtered candidates, and returns the final selected indices in ascending axial order.

If MiniBatchKMeans fitting fails, or if there are too few remaining candidates for clustering, the implementation falls back to evenly spaced indices from the available filtered candidate set. The same uniqueness, validity, and sorting step is then applied before returning \(\mathcal{S}_i\).


\subsection{Final Selection Configuration}
\label{sec:final_selection_configuration}

Table~\ref{tab:slice_selection_configuration} summarizes the final slice-selection configuration. These settings define the slice-selection stage only, where input representation, CNN feature extraction, and patient-level aggregation are described later.

\begin{table}[H]
\centering
\small
\caption{Final configuration for slice selection strategy.}
\label{tab:slice_selection_configuration}
\begin{tabular}{p{0.27\textwidth}p{0.18\textwidth}p{0.45\textwidth}}
\toprule
Parameter & Value & Usage \\
\midrule
\texttt{CENTER\_FRACTION} & \(0.75\) & Retains the central portion of the axial stack before filtering, reducing the influence of extreme scan ends. \\
\texttt{N\_SLICES\_REP} & \(160\) & Requested number of representative slices selected per patient when enough valid candidates are available. \\
\texttt{FEAT\_HW} & \(64\) & Defines the low-resolution spatial size used to construct flattened descriptors for diversity sampling. \\
\texttt{N\_PCA} & \(24\) & Maximum number of PCA components used to compress slice-selection descriptors before clustering. \\
\texttt{LUNG\_HU\_THR} & \(-350\) HU & Threshold used to form a rough lung/air proxy for slice-level filtering and descriptor statistics. \\
\texttt{MIN\_LUNG\_FRAC} & \(0.06\) & Minimum required fraction of pixels below \texttt{LUNG\_HU\_THR}; removes slices with little lung-like content. \\
\texttt{MAX\_AIR\_FRAC} & \(0.98\) & Maximum allowed fraction of very low-density air-like pixels, removes slices that are almost entirely air. \\
\texttt{HI\_DENS\_THR} & \(-200\) HU & Threshold used within the rough lung mask to rank slices by high-density content for informative-slice reservation. \\
\texttt{N\_INFORMATIVE} & \(30\) & Upper cap on the number of high-informativeness slices considered for reserved inclusion before diversity sampling. \\
\texttt{MAX\_BODY\_FRAC} & \(0.45\) & Maximum allowed body/soft-tissue-dominant fraction, reduces inclusion of slices dominated by non-lung tissue. \\
\texttt{HU\_CANON\_MIN} & \(-1350\) HU & Lower bound of the canonical HU-like clipping range used for QC and heuristic slice-selection metrics. \\
\texttt{HU\_CANON\_MAX} & \(600\) HU & Upper bound of the canonical HU-like clipping range used for QC and heuristic slice-selection metrics. \\
\bottomrule
\end{tabular}
\end{table}

The selected indices are returned in ascending axial order before the corresponding slices are exposed to the later input-representation stage. If fewer than \texttt{N\_SLICES\_REP} valid candidates are available after filtering and fallback handling, the implementation returns the available valid subset rather than duplicating slices to force a fixed count.



\section{Input Representation}
\label{sec:input_representation}

After slice selection, each selected preprocessed slice \(\tilde{x}_{i,j}\), \(j \in \mathcal{S}_i\), is converted into a multi-channel array before CNN feature extraction. This stage defines the representation function
\[
r_{i,j}
=
R(\tilde{x}_{i,j}),
\qquad
j \in \mathcal{S}_i,
\]
where \(r_{i,j}\) is the model input representation for the selected slice. In the implemented workflow, \(R(\cdot)\) is either a two-window CT representation or an extended three-channel representation. All channels are stored in channel-first format and resized to the common in-plane size \(H_0=W_0=256\).


\subsection{Baseline Representation (Two-Window)}
\label{sec:baseline_representation_two_window}

The baseline input representation stacks two CT intensity windows: a lung window and a mediastinal window \cite{whiting2015ctbasic,lu2020multiwindow}.

For a window \(w\) with lower and upper HU bounds \(a_w\) and \(b_w\), the implemented windowing operation is
\[
W_w(\tilde{x})
=
\frac{
\min\!\left(\max\!\left(\tilde{x}, a_w\right), b_w\right)-a_w
}{
b_w-a_w+\epsilon
},
\]
with \(\epsilon=10^{-6}\) in the implementation, followed by slice-wise resizing to \(256 \times 256\) and clipping to \([0,1]\). The lung-window bounds are
\[
a_{\mathrm{lung}}=-1350,
\qquad
b_{\mathrm{lung}}=150,
\]
and the mediastinal-window bounds are
\[
a_{\mathrm{med}}=-160,
\qquad
b_{\mathrm{med}}=240.
\]
Thus, for a selected slice \(\tilde{x}_{i,j}\), the two windowed channels are
\[
x_{i,j}^{(\mathrm{lung})}
=
W_{\mathrm{lung}}(\tilde{x}_{i,j}),
\qquad
x_{i,j}^{(\mathrm{med})}
=
W_{\mathrm{med}}(\tilde{x}_{i,j}).
\]
The two-channel baseline representation is then
\[
r_{i,j}^{(2)}
=
\left[
x_{i,j}^{(\mathrm{lung})};
x_{i,j}^{(\mathrm{med})}
\right]
\in
\mathbb{R}^{2 \times 256 \times 256}.
\]
This representation is used as the baseline input format for the two-window experiments and as the shared foundation for the three-channel variants.


\subsection{Extended Representations (Three-Channel Variants)}
\label{sec:extended_representations_three_channel}
The extended representations keep the lung and mediastinal channels and append a third channel derived from a fixed high-frequency or local-structure transform. The high-pass, Difference-of-Gaussians, Laplacian-of-Gaussian, and DCT-based variants are standard image-processing operations, in this thesis they are used only as implemented third-channel construction rules rather than as trainable model components \cite{gonzalez2018digital}. For a third-channel method \(m\), the implemented representation is written as
\[
r_{i,j}^{(3,m)}
=
\left[
x_{i,j}^{(\mathrm{lung})};
x_{i,j}^{(\mathrm{med})};
h_{i,j}^{(m)}
\right],
\qquad
h_{i,j}^{(m)}
=
H_m(u_{i,j}),
\]
where \(u_{i,j}\) is the normalized base channel used to construct the third channel. The implementation supports deriving \(u_{i,j}\) from either the lung-window channel or the mediastinal-window channel through the \texttt{HF\_FROM} setting. In the main high-pass configuration, this source is the mediastinal-window channel.

The main high-frequency variant uses Gaussian blur subtraction. A smoothed version of the base image is subtracted from the original base image, so the resulting response emphasizes local intensity changes relative to the surrounding low-frequency content. With \(G_{\sigma}(\cdot)\) denoting Gaussian smoothing, the implemented high-pass channel is
\[
H_{\mathrm{HP}}(u)
=
N_{\mathrm{HF}}\!\left(u-G_{\sigma_{\mathrm{HP}}}(u)\right),
\qquad
\sigma_{\mathrm{HP}}=2.0,
\]
where \(N_{\mathrm{HF}}(\cdot)\) denotes robust third-channel normalization. In the implementation, this normalization clips the response using the 1st and 99th percentiles and rescales it to \([0,1]\).

The Difference-of-Gaussians variant is implemented as
\[
H_{\mathrm{DoG}}(u)
=
N_{\mathrm{HF}}\!\left(
G_{\sigma_1}(u)-G_{\sigma_2}(u)
\right),
\qquad
\sigma_1=1.0,\quad
\sigma_2=3.0.
\]
This variant compares smoothed versions of the same base channel at two spatial scales and keeps the normalized difference as the third channel.

The Laplacian-of-Gaussian variant first smooths the base channel and then applies a Laplacian operation:
\[
H_{\mathrm{LoG}}(u)
=
N_{\mathrm{HF}}\!\left(
\nabla^2 G_{\sigma_{\mathrm{LoG}}}(u)
\right).
\]
The LoG smoothing scale follows the corresponding run configuration and is not treated as a fixed setting in the main high-pass configuration.

A DCT-based third-channel variant is also implemented. In this variant, DCT coefficients are computed from the base channel, the lowest-frequency \(16 \times 16\) coefficient block is removed, the image is reconstructed with the inverse transform, the absolute value is taken, and the same robust normalization is applied. In the DCT configuration, no additional high-frequency cutoff is applied after the low-frequency removal.

When lung-mask suppression is enabled, the third-channel response is multiplied by a soft rough lung mask. The main high-pass configuration enables this option. The mask suppresses responses outside the approximate lung region, it is not a supervised segmentation target or localization output.



\subsection{Representation Used in Experiments}
\label{sec:representation_used_experiments}

The experiments use the two-window representation and selected three-channel extensions as input formats. The selected representation is summarized as
\[
r_{i,j}
=
\begin{cases}
r_{i,j}^{(2)}, & \text{two-window experiments},\\
r_{i,j}^{(3,m)}, & \text{three-channel experiments with variant } m.
\end{cases}
\]
The two-window baseline uses \(C=2\) input channels. The three-channel experiments use \(C=3\), where the third channel is configured by \texttt{HF\_MODE}. The high-pass representation is the main three-channel configuration, while DoG, LoG, and DCT-based inputs are included as comparison variants.


\begin{table}[H]
\centering
\small
\caption{Main settings for the high-pass input representation.}
\label{tab:input_representation_settings}
\begin{tabular}{p{0.35\textwidth}p{0.22\textwidth}p{0.33\textwidth}}
\toprule
Setting & Value & Use \\
\midrule
Image size & \(256 \times 256\) & Fixed in-plane size for all input channels. \\
Lung window & \([-1350,150]\) HU & First channel in all two-window and three-channel inputs. \\
Mediastinal window & \([-160,240]\) HU & Second channel in all two-window and three-channel inputs. \\
Main high-pass source & \texttt{medi} & Base channel for the main high-pass representation. \\
Main high-pass scale & \(\sigma_{\mathrm{HP}}=2.0\) & Gaussian scale used before high-pass subtraction. \\
HF normalization & 1st--99th percentile & Robust rescaling of third-channel responses to \([0,1]\). \\
HF lung-mask suppression & Enabled for main high-pass & Suppresses third-channel response outside the approximate lung region. \\
\bottomrule
\end{tabular}
\end{table}



Figure~\ref{fig:highpass_masking_example} shows the high-pass third channel before and after mediastinal-mask suppression.

\begin{figure}[H]
\centering
\begin{subfigure}[t]{0.38\textwidth}
    \centering
    \includegraphics[width=\linewidth]{figures/highpass_unmasked_example.png}
    \caption{Without suppression.}
    \label{fig:highpass_unmasked}
\end{subfigure}
\hspace{0.04\textwidth}
\begin{subfigure}[t]{0.38\textwidth}
    \centering
    \includegraphics[width=\linewidth]{figures/highpass_masked_example.png}
    \caption{With suppression.}
    \label{fig:highpass_masked}
\end{subfigure}
\caption{High-pass third-channel representation before and after lung-mask suppression.}
\label{fig:highpass_masking_example}
\end{figure}

The representation stage ends with a channel-first tensor \(r_{i,j}\), which is passed to the CNN feature-extraction stage in Section~\ref{sec:feature_extraction_model_architecture}.


\section{Feature Extraction and Model Architecture}
\label{sec:feature_extraction_model_architecture}

This section describes the slice-level mapping from input representation \(r_{i,j}\) to learned feature vector \(z_{i,j}\), slice score \(s_{i,j}\), and in fusion experiments the combined learned--handcrafted representation. For a selected slice \(j \in \mathcal{S}_i\), the standard CNN path is
\[
z_{i,j}
=
\Phi_{\psi}(r_{i,j}),
\qquad
s_{i,j}
=
g_{\omega}(z_{i,j}),
\]
where \(\Phi_{\psi}\) denotes the CNN feature extractor with parameters \(\psi\), and \(g_{\omega}\) denotes the trainable regression head with parameters \(\omega\).

When handcrafted features are enabled, the learned CNN feature vector is combined with a handcrafted feature vector before slice-level prediction. Let \(b_{i,j}\) denote the original handcrafted feature vector for slice \(j\), and let \(q_{\eta}\) denote the trainable handcrafted projection branch with parameters \(\eta\). The fusion path is written as
\[
\bar{b}_{i,j}
=
q_{\eta}(b_{i,j}),
\qquad
v_{i,j}
=
\left[z_{i,j}; \bar{b}_{i,j}\right],
\qquad
s_{i,j}
=
g_{\omega}(v_{i,j}).
\]
Here \(\bar{b}_{i,j}\) is the projected handcrafted representation and \(v_{i,j}\) is the fused feature vector. Patient-level aggregation is described in Section~\ref{sec:patient_level_prediction_aggregation}.


\subsection{CNN Backbones}
\label{sec:cnn_backbones}

The CNN input for each selected slice is the channel-first tensor
\[
r_{i,j}
\in
\mathbb{R}^{C \times 256 \times 256},
\qquad
j \in \mathcal{S}_i,
\]
where \(C=2\) for the two-window representation and \(C=3\) for the three-channel variants described in Section~\ref{sec:input_representation}. The implementation defines a backbone-dependent feature extractor
\[
z_{i,j}^{(B)}
=
\Phi_{\psi}^{(B)}(r_{i,j}),
\]
where \(B\) denotes the selected backbone type. The corresponding slice-level prediction is
\[
s_{i,j}^{(B)}
=
g_{\omega}^{(B)}(z_{i,j}^{(B)}).
\]

The implementation uses torchvision backbones as feature extractors. Since these backbones expect three-channel input, the model includes an input projection stage. If \(C=3\), this projection is the identity map. If \(C \neq 3\), then add a trainable \(1 \times 1\) convolution from \(C\) channels to three channels:
\[
r'_{i,j}
=
P_{\mathrm{in}}(r_{i,j}),
\qquad
P_{\mathrm{in}}:
\mathbb{R}^{C \times 256 \times 256}
\rightarrow
\mathbb{R}^{3 \times 256 \times 256}.
\]



\subsection{ResNet18 Architecture}
\label{sec:resnet18_architecture}

The main backbone configuration uses torchvision ResNet18 \cite{he2016residual}. The ResNet18 classification layer is replaced with \texttt{nn.Identity()}, so the network outputs a feature vector rather than class logits. For a selected slice, the feature path is
\[
z_{i,j}^{(\mathrm{ResNet})}
=
\Phi_{\psi}^{(\mathrm{ResNet18})}(r_{i,j}),
\qquad
z_{i,j}^{(\mathrm{ResNet})}
\in
\mathbb{R}^{512}.
\]
In the main three-channel configuration, \(C=3\), so the pretrained ResNet18 input stem is kept unchanged. In the two-window configuration, a trainable \(1 \times 1\) convolution maps the two channels to the expected three-channel input.

The slice-level regression head is a small multilayer head applied after layer normalization:
\[
s_{i,j}^{(\mathrm{ResNet})}
=
g_{\omega}^{(\mathrm{ResNet18})}
\!\left(z_{i,j}^{(\mathrm{ResNet})}\right),
\qquad
g_{\omega}^{(\mathrm{ResNet18})}:
\mathbb{R}^{512}
\rightarrow
\mathbb{R}.
\]
The head consists of a linear layer to 128 units, ReLU activation, dropout, and a final linear layer to one scalar output. The main ResNet18 configuration uses torchvision's default pretrained weights when no RadImageNet checkpoint path is supplied.


\subsection{DenseNet121 Architecture}
\label{sec:densenet121_architecture}

DenseNet121 is implemented as a comparison backbone using torchvision DenseNet121 \cite{huang2017densenet}. Its classification layer is replaced by an identity mapping. For a selected slice,
\[
z_{i,j}^{(\mathrm{DenseNet})}
=
\Phi_{\psi}^{(\mathrm{DenseNet121})}(r_{i,j}),
\qquad
z_{i,j}^{(\mathrm{DenseNet})}
\in
\mathbb{R}^{1024}.
\]
As with ResNet18, the input representation is first mapped to three channels if needed. The DenseNet121 slice-level score is produced by the same type of normalized MLP regression head:
\[
s_{i,j}^{(\mathrm{DenseNet})}
=
g_{\omega}^{(\mathrm{DenseNet121})}
\!\left(z_{i,j}^{(\mathrm{DenseNet})}\right),
\qquad
g_{\omega}^{(\mathrm{DenseNet121})}:
\mathbb{R}^{1024}
\rightarrow
\mathbb{R}.
\]
The implementation supports ImageNet initialization through torchvision and external RadImageNet weights. In the reported comparison, DenseNet121 uses a RadImageNet checkpoint path, whereas ResNet18 uses ImageNet initialization.


\subsection{Handcrafted Feature Extraction}
\label{sec:handcrafted_feature_extraction}

Some experiments supplement the learned CNN feature vector with a handcrafted feature vector. This branch is separate from the slice-selection descriptors \(d_{i,j}\) and \(p_{i,j}\) used earlier; those descriptors choose representative slices, while the handcrafted features here are passed to the regression model.

For a selected slice, the handcrafted feature vector is written as
\[
b_{i,j}
=
B_{\mathrm{BoVW}}(\tilde{x}_{i,j}),
\qquad
b_{i,j}
\in
\mathbb{R}^{K_{\mathrm{vw}}},
\]
where \(B_{\mathrm{BoVW}}(\cdot)\) denotes the ORB-based bag-of-visual-words extraction procedure and \(K_{\mathrm{vw}}\) is the visual vocabulary size. In the implemented fusion configuration,
\[
K_{\mathrm{vw}} = 64.
\]
The handcrafted branch is a fixed feature-extraction step followed by a trainable projection inside the model.


\subsection{ORB and Bag-of-Visual-Words}
\label{sec:orb_bovw}

The handcrafted branch uses ORB descriptors computed with OpenCV's ORB implementation \cite{rublee2011orb}. In the main fusion configuration, the keypoint image is the mediastinal-window image, with \texttt{KEYPOINT\_SOURCE}=\texttt{medi}, and keypoint detection is restricted by the mediastinal mask used by the dataset helper.

Let
\[
\mathcal{D}_{i,j}
=
\{\delta_{i,j,1},\dots,\delta_{i,j,m_{i,j}}\}
\]
denote the variable-length set of ORB descriptors extracted for slice \(j\) from patient \(i\). The visual vocabulary is fitted on training-fold descriptors only, using MiniBatchKMeans, so that validation slices are encoded using a vocabulary not fitted on their descriptors. The vocabulary is written as
\[
\mathcal{V}
=
\{c_1,\dots,c_{K_{\mathrm{vw}}}\},
\qquad
K_{\mathrm{vw}}=64.
\]
The vocabulary configuration uses at most 50,000 descriptors for fitting, 500 ORB features per slice, MiniBatchKMeans batch size 4096, maximum 200 iterations, and a fold-dependent random seed.

Each slice descriptor set is converted into a normalized visual-word histogram. For visual word \(k\), the histogram entry can be written as
\[
b_{i,j,k}
=
\frac{1}{m_{i,j}}
\sum_{\ell=1}^{m_{i,j}}
\mathbf{1}
\!\left[
k
=
\arg\min_{1 \leq q \leq K_{\mathrm{vw}}}
\|\delta_{i,j,\ell}-c_q\|_2
\right],
\]
when at least one descriptor is available. The resulting vector is
\[
b_{i,j}
=
\left[
b_{i,j,1},\dots,b_{i,j,K_{\mathrm{vw}}}
\right]
\in
\mathbb{R}^{K_{\mathrm{vw}}}.
\]
This follows the standard bag-of-visual-words idea of representing variable-length local descriptors by a fixed-length visual-word histogram, which has also been used in CT image retrieval settings \cite{yang2012bovwct}. If ORB finds no descriptors for a slice, the implementation returns the zero vector in \(\mathbb{R}^{K_{\mathrm{vw}}}\).


\subsection{Feature Fusion Strategy}
\label{sec:feature_fusion_strategy}

Feature fusion is performed at slice level before patient-level aggregation. When handcrafted fusion is disabled, the regression head operates directly on the CNN feature vector:
\[
s_{i,j}
=
g_{\omega}(z_{i,j}).
\]
When ORB-BoVW fusion is enabled, the BoVW histogram is first passed through a trainable projection branch consisting of layer normalization, a linear projection to 64 dimensions, and ReLU activation:
\[
\bar{b}_{i,j}
=
q_{\eta}(b_{i,j}),
\qquad
\bar{b}_{i,j}
\in
\mathbb{R}^{64},
\]
where \(\eta\) denotes the trainable parameters of the handcrafted projection branch. The projected handcrafted vector is then concatenated with the normalized CNN feature vector:
\[
v_{i,j}
=
\left[
z_{i,j};
\bar{b}_{i,j}
\right]
\in
\mathbb{R}^{d_z+64},
\]
where \(d_z=512\) for ResNet18 and \(d_z=1024\) for DenseNet121. The slice-level score is then computed as
\[
s_{i,j}
=
g_{\omega}(v_{i,j}).
\]
The fused regression head has the same output structure as the CNN-only head, a linear layer to 128 units, ReLU activation, dropout, and a final scalar output. Missing cached BoVW features are replaced by zero vectors. The fusion output remains a slice-level score and is aggregated later in Section~\ref{sec:patient_level_prediction_aggregation}.



\section{Patient-Level Prediction Aggregation}
\label{sec:patient_level_prediction_aggregation}

The regression head produces one scalar score for each selected slice, whereas \(y_i\) is patient-level. The pipeline therefore aggregates slice-level scores into one patient-level prediction:
\[
\hat{y}_i
=
A\left(\{s_{i,j}: j \in \mathcal{S}_i\}\right),
\]
where \(A(\cdot)\) denotes the selected patient-level aggregation rule. The implementation supports mean aggregation, top-\(k\) aggregation, and a hybrid aggregation rule that combines the two.


\subsection{Slice-Level Predictions}
\label{sec:slice_level_predictions}

For each selected slice \(j \in \mathcal{S}_i\), the model computes a slice-level score from the learned feature vector:
\[
s_{i,j}
=
g_{\omega}(z_{i,j}).
\]
When handcrafted feature fusion is enabled, the same prediction head is applied to the fused feature vector:
\[
s_{i,j}
=
g_{\omega}(v_{i,j}).
\]
For patient \(i\), the scores available for aggregation are collected as
\[
\mathbf{s}_i
=
\{s_{i,j}: j \in \mathcal{S}_i\},
\qquad
M_i = |\mathcal{S}_i|.
\]
The slice-level scores are intermediate outputs, not independently validated slice-level targets. If an experiment is trained on a transformed target, aggregation is applied in the corresponding model-output scale. The implementation includes an empty-list fallback, but slice selection is designed to provide a non-empty selected set for each case.


\subsection{Mean Aggregation}
\label{sec:mean_aggregation}

Mean aggregation averages all selected slice-level scores:
\[
\hat{y}_i^{(\mathrm{mean})}
=
\frac{1}{M_i}
\sum_{j \in \mathcal{S}_i}
s_{i,j},
\qquad
M_i = |\mathcal{S}_i|.
\]
This mode uses every selected slice score with equal weight and has no additional hyperparameter.


\subsection[Top-k Aggregation]{Top-\(k\) Aggregation}
\label{sec:topk_aggregation}

Top-\(k\) aggregation uses the largest predicted slice scores for a patient. The implementation supports either a fixed integer \(k\) or a fraction of the selected slices. When fixed \(k\) is supplied, it takes priority. The effective count is capped by the available slice predictions:
\[
k_i^{(\mathrm{eff})}
=
\min(k, M_i).
\]
Let \(\mathcal{T}_{i,k} \subseteq \mathcal{S}_i\) denote the indices of the \(k_i^{(\mathrm{eff})}\) largest slice scores for patient \(i\):
\[
\mathcal{T}_{i,k}
=
\operatorname{TopK}_{k_i^{(\mathrm{eff})}}
\left(\{s_{i,j}: j \in \mathcal{S}_i\}\right).
\]
The patient-level prediction under top-\(k\) aggregation is then
\[
\hat{y}_i^{(\mathrm{top}k)}
=
\frac{1}{k_i^{(\mathrm{eff})}}
\sum_{j \in \mathcal{T}_{i,k}}
s_{i,j}.
\]
In the main aggregation configuration, the fixed setting is \(k=20\). The implementation also keeps a fractional top-\(k\) option. If a fractional setting is used instead, the effective count is computed from \(M_i\) and constrained to be at least one.


\subsection{Hybrid Aggregation Strategy}
\label{sec:hybrid_aggregation_strategy}

The hybrid aggregation mode, named \texttt{mix} in the implementation, combines the top-\(k\) and mean estimates:
\[
\hat{y}_i^{(\mathrm{hybrid})}
=
\lambda
\hat{y}_i^{(\mathrm{top}k)}
+
(1-\lambda)
\hat{y}_i^{(\mathrm{mean})}.
\]
The mixing weight \(\lambda\) is a fixed aggregation setting, not a learned model parameter. In the main configuration,
\[
\lambda = 0.7,
\qquad
k = 20.
\]
Thus, the hybrid rule gives greater weight to the top-\(k\) component while retaining a contribution from all selected slices. The configured fractional value \(\texttt{TOPK\_FRAC}=0.15\) is present, but it is not used when fixed \(\texttt{TOPK\_K}=20\) is supplied.


\section{Training, Transfer Learning, and Model Selection}
\label{sec:training_transfer_model_selection}

This section describes the training objective, optimization settings, fold protocol, transfer-learning setup, and model-selection rule. Quantitative validation results are reported in Chapter~5.

\subsection{Learning Objective}
\label{sec:learning_objective}

The final prediction target is defined at patient level, but the implemented training loop optimizes sampled slice predictions using the corresponding patient label as a weak training signal. In this section, \(y_i\) denotes the selected target variable for patient \(i\). For the main internal training runs this target is the log-transformed label field \texttt{label\_log10}, the same training code can also be applied to another selected target column.

For a mini-batch of sampled slices \(\mathcal{B}_{\mathrm{slice}}\), where each element is a patient--slice pair \((i,j)\), the model computes a slice-level score \(s_{i,j}\). The implemented loss is mean squared error:
\begin{equation}
\mathcal{L}_{\mathrm{MSE}}
=
\frac{1}{|\mathcal{B}_{\mathrm{slice}}|}
\sum_{(i,j)\in\mathcal{B}_{\mathrm{slice}}}
\left(y_i - s_{i,j}\right)^2 .
\label{eq:method_mse_slice_training}
\end{equation}
The label used in the loss is patient-level, but optimization is carried out over sampled slice predictions. Patient-level predictions
\[
\hat{y}_i =
A\left(\{s_{i,j}:j\in\mathcal{S}_i\}\right)
\]
are computed during validation and model selection using the aggregation rules described in Section~\ref{sec:patient_level_prediction_aggregation}.


\subsection{Optimization Settings}
\label{sec:optimization_settings}
The trainable parameters depend on whether the handcrafted branch is enabled. Without fusion the optimized parameter set consists of the CNN feature extractor and prediction head
\[
\Theta=\{\psi,\omega\}.
\]
When ORB/BoVW fusion is enabled the handcrafted projection branch is also trained
\[
\Theta=\{\psi,\eta,\omega\}.
\]
Here \(\eta\) denotes the handcrafted-feature projection branch. Aggregation settings such as \(k\) and \(\lambda\) are fixed, not trainable.

The main internal training configuration uses AdamW optimization, mean squared error loss, and early stopping based on validation performance. The values below are taken from the implemented high-pass and ORB/BoVW main training configuration.

\begin{table}[htbp]
\centering
\caption{Main training settings used in the implemented internal cross-validation runs.}
\label{tab:method_training_settings}
\begin{tabular}{p{0.27\textwidth}p{0.22\textwidth}p{0.41\textwidth}}
\toprule
Setting & Value & Note \\
\midrule
Optimizer & AdamW & Applied to parameters with \texttt{requires\_grad=True} \\
Initial learning rate & \(1\times10^{-4}\) & Used before partial unfreezing \\
Learning rate after unfreezing & \(3\times10^{-5}\) & Implemented as \(0.3\) times the initial rate \\
Weight decay & \(2\times10^{-4}\) & AdamW weight decay \\
Batch size & \(32\) & Slice-level mini-batches \\
Maximum epochs & \(30\) & Per fold \\
Early-stopping patience & \(6\) epochs & Based on validation RMSE in log space \\
Dropout & \(0.3\) & Applied in the prediction head configuration \\
Gradient clipping & \(1.0\) & Norm clipping when enabled \\
Data augmentation & Enabled & Rotation up to \(5^\circ\), translation up to \(6\) pixels, intensity jitter \(0.03\), and noise \(0.02\) \\
\bottomrule
\end{tabular}
\end{table}

The main training loop does not use a separate learning-rate scheduler. Instead, the optimizer is recreated with a lower learning rate when later backbone layers are unfrozen.

\subsection{Cross-Validation Protocol}
\label{sec:cross_validation_protocol}
Model development is organized as five-fold cross-validation using fold assignments stored in the locked split file. For fold \(f\), the validation set contains patients assigned to that fold, while the training set contains the remaining patients:
\begin{equation}
\mathcal{I}
=
\mathcal{I}_{\mathrm{train}}^{(f)}
\cup
\mathcal{I}_{\mathrm{val}}^{(f)},
\qquad
\mathcal{I}_{\mathrm{train}}^{(f)}
\cap
\mathcal{I}_{\mathrm{val}}^{(f)}
=
\emptyset .
\label{eq:method_cv_split}
\end{equation}
The split is patient-level, so slices from the same patient are not divided between training and validation.

For feature-fusion experiments, the BoVW vocabulary is fitted only on the training part of each fold. Validation slices are encoded using the corresponding training vocabulary, keeping the handcrafted branch within the cross-validation boundary.


\subsection{Transfer Learning Setup}
\label{sec:method_transfer_learning_setup}
Transfer learning appears in two forms: pretrained image-model initialization for internal cross-validation, and staged fine-tuning from a source-domain checkpoint for the external source--target transfer workflow.


\subsubsection{Pretraining Sources}
\label{sec:pretraining_sources}
For the main ResNet18 configuration, the backbone is initialized with ImageNet weights from \texttt{torchvision}. The DenseNet121 comparison configuration supports initialization from a local RadImageNet checkpoint, a radiology-domain pretraining resource \cite{mei2022radimagenet}. ResNet18 and DenseNet121 are the relevant backbones for this methodology chapter.

In the external transfer workflow, the source initialization is a checkpoint trained on the corresponding source dataset. Mucus-to-MosMed transfer starts from an internal mucus-score checkpoint, while MosMed-to-mucus transfer starts from a MosMed severity checkpoint. This is separate from ImageNet or RadImageNet initialization because it transfers from a task-specific CT model.


\subsubsection{Model Initialization}
\label{sec:model_initialization}
Pretrained weights initialize the backbone feature extractor:
\begin{equation}
\psi_0 \leftarrow \psi_{\mathrm{pre}} .
\label{eq:method_pretrained_backbone_init}
\end{equation}
The original classification layer is replaced by a newly initialized regression head:
\begin{equation}
\omega_0 \sim \mathrm{Init}(\cdot).
\label{eq:method_head_init}
\end{equation}
When feature fusion is enabled the handcrafted projection branch is also initialized newly:
\begin{equation}
\eta_0 \sim \mathrm{Init}(\cdot).
\label{eq:method_fusion_init}
\end{equation}

For three-channel inputs, the represented slice matches the pretrained backbone input format. For two-channel inputs, a newly initialized projection layer maps the representation to three channels and is trained with the non-frozen model components.


\subsubsection{Fine-Tuning Strategy}
\label{sec:fine_tuning_strategy}
Internal cross-validation uses staged fine-tuning. Training begins with the pretrained backbone frozen, so only newly added parts such as the regression head, input projection layer, and handcrafted projection branch are optimized. After the frozen phase, the final backbone stage is unfrozen and optimized with a smaller learning rate.

In the main configuration, the frozen phase lasts for \(2\) epochs. After that, ResNet18 unfreezes \texttt{layer4}, while DenseNet121 unfreezes \texttt{features.denseblock4} and \texttt{features.norm5}. The optimizer is recreated with learning rate \(3\times10^{-5}\), so adaptation is limited to later backbone layers.

The external transfer workflow follows the same idea, it starts from the source-domain checkpoint, trains the head-only configuration on the target dataset, then unfreezes \texttt{layer4}, with optional \texttt{layer3} unfreezing for a final low-learning-rate stage.


\subsubsection{External Source--Target Transfer}
\label{sec:external_source_target_transfer}
The external transfer experiments use a source-domain checkpoint trained on the corresponding source dataset and evaluate reuse on a different target dataset. Mucus-to-MosMed transfer uses an internal mucus-score checkpoint and adapts it to MosMed. MosMed-to-mucus transfer uses a MosMed severity checkpoint and adapts it to the internal mucus-score task. The purpose is cross-domain reuse of CT-based representations, not external validation for mucus-burden prediction.

For each direction, zero-shot evaluation applies the source-domain checkpoint to the target domain without updating model parameters. Fine-tuning uses the same checkpoint as initialization and trains further on target-domain training data before validation.

Because the MosMed and mucus tasks use different targets and output scales, MosMed-to-mucus transfer applies a fixed scale mapping when MosMed-scale predictions are compared with mucus-score targets:
\[
\hat{y}_{i}^{(\mathrm{raw})}
=
2\,\hat{y}_{i}^{(\mathrm{M})}.
\]
This is a fixed implementation choice and should not be interpreted as a learned calibration between the two clinical targets.

When MosMed-to-mucus fine-tuning is performed against the log-transformed mucus target, the scaled raw prediction is clipped and transformed before computing the loss:
\[
\hat{y}_{i}^{(\log)}
=
\log_{10}
\left(
\operatorname{clip}
\left(
2\,\hat{y}_{i}^{(\mathrm{M})},
10^{-6},
8.0
\right)
\right).
\]
The model is then optimized against the stored mucus \texttt{label\_log10} target. For reporting on the mucus task, predictions are evaluated on the raw mucus-score scale when raw-scale metrics are shown.

Transfer results are computed at case level after slice-level predictions have been aggregated. Chapter~5 reports the mean metric value across folds for each transfer setting.


\subsection{Model Selection Criteria}
\label{sec:model_selection_criteria}
Model selection is performed separately for each fold using patient-level validation predictions. After each epoch \(t\), validation slice scores are aggregated into one prediction \(\hat{y}_i^{(t)}\) for each validation patient \(i \in \mathcal{I}_{\mathrm{val}}^{(f)}\). The checkpoint is selected using validation RMSE in the same target space used for training. For the main internal experiments, this is the log-transformed target space,
\[
\mathrm{RMSE}_{\mathrm{val}}^{(f)}(t)
=
\sqrt{
\frac{1}{|\mathcal{I}_{\mathrm{val}}^{(f)}|}
\sum_{i\in\mathcal{I}_{\mathrm{val}}^{(f)}}
\left(y_i-\hat{y}_i^{(t)}\right)^2
}.
\]
Here, \(y_i\) denotes the selected training target for the experiment, such as the log-transformed mucus target in the main internal runs. The selected epoch for fold \(f\) is therefore
\[
t_f^{*}
=
\arg\min_t
\mathrm{RMSE}_{\mathrm{val}}^{(f)}(t).
\]
The best model state for each fold is stored during training and restored after early stopping. MAE, \(R^2\), and Spearman correlation are also computed for reporting, but they are not the primary checkpoint-selection criterion.

For experiments trained on the log-transformed target, model selection is done in log space, while the final reported MAE and RMSE can be computed after transforming predictions back to the raw mucus-score scale. This keeps checkpoint selection consistent with the training objective, while keeping the reported results easier to interpret.


\section{Implementation and Reproducibility}
\label{sec:implementation_reproducibility}
This section records the practical setup used to run the experiments. The purpose is not to repeat the full pipeline, but to document the software stack, main hyperparameters, randomness controls, and computational setup needed to reproduce the methodology as closely as possible.


\subsection{Software and Libraries}
\label{sec:software_libraries}
The experiments were implemented in Python notebooks using Google Colab Pro and Google Drive. The software documentation versions are listed in Table~\ref{tab:software_libraries}. 

\begin{table}[H]
\centering
\small
\caption{Main software libraries used in the implementation.}
\label{tab:software_libraries}
\begin{tabular}{p{0.25\textwidth}p{0.20\textwidth}p{0.45\textwidth}}
\toprule
Library or tool & Documented version & Role in the implementation \\
\midrule
Python & 3.12.13 & Notebook execution and experiment orchestration. \\
PyTorch & 2.10.0 & Tensor operations, CNN model definition, loss computation, optimization, and checkpoint handling. \\
\texttt{torchvision} & 0.25.0 & Pretrained CNN backbones, including ResNet18 and DenseNet121. \\
NumPy & 2.0.2 & Stored-volume loading, array operations, slice descriptors, and selected random generators. \\
pandas & 2.2.2 & Dataset indexing, fold assignment, label fields, and tabular summaries. \\
OpenCV & 4.13.0 & Slice resizing, CT channel construction, filtering, DCT transforms, mask operations, and ORB descriptors. \\
scikit-learn & 1.6.1 & PCA, MiniBatchKMeans, nearest-centroid assignment, and BoVW vocabulary construction. \\
SciPy & 1.16.3 & Correlation computation during validation reporting. \\
Matplotlib & 3.10.0 & Quality-control plots and visual inspection figures. \\
nibabel & 5.4.2 & Loading NIfTI volumes in the external MosMed transfer workflow. \\
Google Colab Pro & -- & Notebook runtime environment and GPU access. \\
\bottomrule
\end{tabular}
\end{table}


\subsection{Hyperparameters}
\label{sec:hyperparameters}
Stage-specific settings are described in the relevant methodology sections. Tables~\ref{tab:main_pipeline_hyperparameters} and~\ref{tab:training_repro_hyperparameters} collect the main confirmed values for the \textbf{ResNet18-3CH+BoVW (Proposed model)} configuration.

\begin{table}[H]
\centering
\small
\caption{Main data, slice-selection, and representation hyperparameters.}
\label{tab:main_pipeline_hyperparameters}
\begin{tabular}{p{0.24\textwidth}p{0.30\textwidth}p{0.36\textwidth}}
\toprule
Stage & Hyperparameter & Confirmed value \\
\midrule
Preprocessing & Canonical HU clipping & \([-1350, 600]\) HU \\
Input representation & Final in-plane size & \(256 \times 256\) pixels \\
Input representation & Lung window & \([-1350, 150]\) HU \\
Input representation & Mediastinal window & \([-160, 240]\) HU \\
Third channel & Main high-pass source & Mediastinal-window channel \\
Third channel & High-pass Gaussian scale & \(\sigma_{\mathrm{HP}}=2.0\) \\
Third channel & Robust response scaling & 1st--99th percentile normalization \\
Third channel & Lung-mask suppression & Enabled in the main high-pass configuration \\
Slice selection & Central fraction & \(0.75\) \\
Slice selection & Selected slices per case & \(160\) \\
Slice selection & Low-resolution descriptor size & \(64 \times 64\) \\
Slice selection & PCA components & \(24\) \\
Slice selection & Informative-slice cap & \(N_{\mathrm{informative}}=30\) \\
Slice filtering & Main QC thresholds & \(\tau_{\mathrm{lung}}=0.06\), \(\tau_{\mathrm{air}}=0.98\), \(\tau_{\mathrm{body}}=0.45\), high-density threshold \(-200\) HU \\
Aggregation & Main aggregation mode & Hybrid \texttt{mix} aggregation \\
Aggregation & Top-\(k\) setting & \(k=20\), with fractional fallback \(0.15\) \\
Aggregation & Hybrid mixing weight & \(\lambda=0.7\) \\
\bottomrule
\end{tabular}
\end{table}

\begin{table}[H]
\centering
\small
\caption{Main model, fusion, and training hyperparameters.}
\label{tab:training_repro_hyperparameters}
\begin{tabular}{p{0.24\textwidth}p{0.30\textwidth}p{0.36\textwidth}}
\toprule
Stage & Hyperparameter & Confirmed value \\
\midrule
Backbone & Main backbone & ResNet18 \\
Backbone & Input channels & \(C=3\) for the main high-pass configuration \\
Feature fusion & ORB/BoVW branch & Enabled in the main configuration \\
Feature fusion & Visual vocabulary size & \(K_{\mathrm{vw}}=64\) \\
Feature fusion & ORB keypoint limit & \(500\) keypoints per slice \\
Feature fusion & Descriptor cap for vocabulary fitting & \(50{,}000\) descriptors \\
Feature fusion & Keypoint source & Mediastinal-window channel, restricted by the approximate lung mask \\
BoVW clustering & MiniBatchKMeans settings & Batch size \(4096\), \(200\) maximum iterations, \texttt{n\_init=auto} \\
Training & Optimizer & AdamW \\
Training & Batch size & \(32\) \\
Training & Maximum epochs & \(30\) \\
Training & Early-stopping patience & \(6\) epochs \\
Training & Initial learning rate & \(1\times10^{-4}\) \\
Training & Learning rate after partial unfreezing & \(3\times10^{-5}\) \\
Training & Weight decay & \(2\times10^{-4}\) \\
Training & Dropout & \(0.3\) \\
Training & Frozen-backbone phase & \(2\) epochs \\
Training & Gradient clipping & \(1.0\) \\
Data loading & Worker processes & \texttt{num\_workers=0} \\
\bottomrule
\end{tabular}
\end{table}

Comparison representations such as DoG, LoG, and DCT use the same input-representation interface but are not repeated here because they are not the main final configuration.


\subsection{Reproducibility and Randomness Control}
\label{sec:reproducibility_randomness_control}
The implementation uses fixed seeds for selected preprocessing and fold-level operations, including \texttt{RANDOM\_STATE=42} for slice selection, \texttt{SEED=42} in the main training configuration, and fold-specific MiniBatchKMeans seeds defined as \texttt{seed + 1000 + fold}. Fold assignments are stored in a locked split file and reused directly. Within each fold, BoVW vocabularies are fitted only on the training split, with validation slices encoded using the corresponding training vocabulary, preventing validation leakage into the handcrafted features.

The setup is therefore partially controlled but not bitwise deterministic. The notebooks do not define a single global seeding routine for Python \texttt{random}, NumPy, PyTorch, CUDA, and cuDNN, and no explicit cuDNN deterministic settings were confirmed. Training augmentations also include random geometric and intensity transformations. Part of the representative dataset seed construction depends on Python's \texttt{hash} function, which can vary across processes unless \texttt{PYTHONHASHSEED} is fixed. Thus, the split structure and main configuration values are reproducible, while exact slice order and training trajectories may still vary across sessions or hardware backends.


\subsection{Hardware and Computational Resources}
\label{sec:hardware_computational_resources}
Experiments were run in Google Colab Pro, with notebooks configured to use CUDA when available. GPU runtimes were allocated dynamically and could vary between sessions; available options during the project included NVIDIA H100, A100, and G4-based runtimes \cite{nvidia2024h100,nvidia2024a100,googlecloud2025g4gpu}. The documented environment snapshot was CPU-only, reporting \texttt{CUDA available=False} and \texttt{torch.version.cuda=None}, and should therefore be interpreted as package-version documentation rather than a complete record of all GPU hardware used during training.

\chapter{Experiments and Results}
\label{sec:experiments_and_results}

This chapter presents the experimental setup and the quantitative results from the implemented pipeline. The focus is on reporting what was tested and what the models achieved. Deeper interpretation of why certain components helped or failed is kept for Chapter~6.


\section{Experimental Protocol}
\label{sec:experimental_protocol}

All experiments were evaluated at patient level. For each patient \(i\), the model produced one final prediction \(\hat{y}_i\), which was compared with the corresponding patient-level target \(y_i\). Slice-level scores were only used as intermediate outputs inside the model and were not evaluated separately, since no slice-level target labels were available.


\subsection{Internal Five-Fold Cross-Validation}
\label{sec:internal_five_fold_cv_results}

The internal mucus-score experiments used the five-fold cross-validation setup described in Chapter~4. In each fold, the model was trained on the training patients and evaluated on the held-out validation patients. The same patient-level fold assignments were used across the main internal experiments so that different model configurations could be compared under the same validation structure. Results from the internal experiments are reported as mean \(\pm\) standard deviation across the five folds.

Baseline experiments were trained using the raw mucus-score target, while others used the log-transformed target described in Chapter~4. For easier interpretation, the reported error metrics are transformed back to the raw mucus-score scale whenever possible, so that MAE and RMSE can be read in terms of the original burden score.

The internal experiments compare slice handling, input representation, feature fusion, and backbone choice. The two-window ResNet18 model is used as the main reference setting. Additional experiments evaluate the representative slice-selection strategy, three-channel input representations, ORB-BoVW feature fusion, and DenseNet121 as a comparison backbone.


\subsection{Transfer-Learning Evaluation}
\label{sec:transfer_learning_evaluation_protocol}

Transfer-learning experiments are reported separately from the internal five-fold experiments. These experiments evaluate how models behave when trained or initialized across the mucus dataset and the external MosMed dataset. For these experiments, the tables specify which dataset is used as the source domain and which dataset is used as the target domain.

The MosMed experiments are transfer-learning and domain-shift experiments, not external validation of mucus-burden prediction. MosMed uses a different case-level CT severity target, so mucus-to-MosMed evaluates adaptation to the MosMed severity task, while MosMed-to-mucus evaluates adaptation back to the internal mucus-burden task.

The transfer results are reported as average performance across the five folds. The target was modelled on the log-transformed scale, and predictions were transformed back to the raw target scale for reporting. This makes the reported MAE and RMSE easier to interpret in terms of the original score values.

In the transfer tables, zero-shot evaluation refers to applying a source-trained model to the target domain without target-domain training. Fine-tuned evaluation refers to adapting the source-initialized model using target-domain training data before evaluation.



\section{Evaluation Metrics}
\label{sec:results_metrics}

The models are evaluated using patient-level regression metrics. Let \(y_i\) denote the observed target for patient \(i\), and let \(\hat{y}_i\) denote the corresponding patient-level prediction. For a validation set with \(N\) patients, the mean absolute error is defined as
\[
\mathrm{MAE}
=
\frac{1}{N}
\sum_{i=1}^{N}
|y_i-\hat{y}_i|.
\]
MAE measures the average absolute size of the prediction error. Lower values indicate better performance.

The root mean squared error is defined as
\[
\mathrm{RMSE}
=
\sqrt{
\frac{1}{N}
\sum_{i=1}^{N}
(y_i-\hat{y}_i)^2
}.
\]
RMSE also measures prediction error, but it gives larger errors more weight than MAE because the residuals are squared before averaging. Lower RMSE is therefore better. MAE, RMSE, and \(R^2\) are commonly used regression evaluation metrics, but they emphasize different aspects of prediction error and should therefore be interpreted together rather than in isolation \cite{chicco2021coefficient}.

The coefficient of determination is computed as
\[
R^2
=
1 -
\frac{
\sum_{i=1}^{N}(y_i-\hat{y}_i)^2
}{
\sum_{i=1}^{N}(y_i-\bar{y})^2
},
\]
where \(\bar{y}\) is the mean target value in the evaluated set. Higher \(R^2\) is better, with \(1\) representing perfect prediction. However, \(R^2\) should be interpreted relative to the mean-prediction reference. An \(R^2\) value close to zero means that the model explains little additional variation beyond predicting the mean target value, while a negative \(R^2\) means that the model performs worse than this reference for that evaluation set. Since the validation folds are small in this thesis, \(R^2\) is interpreted together with MAE, RMSE, and Spearman correlation rather than as the only summary of performance \cite{chicco2021coefficient}.

Spearman correlation is used to measure rank association between the observed targets and the model predictions \cite{spearman1904proof}. This metric is useful because it shows whether the model tends to order patients correctly by predicted burden, even when the absolute prediction errors are still large. Higher Spearman correlation indicates stronger agreement in ranking.

In addition to reporting mean \(\pm\) standard deviation across the five folds, an approximate confidence interval is reported for the proposed model. For a metric value \(m_f\) computed on fold \(f\), with fold mean \(\bar{m}\), sample standard deviation \(s_m\), and \(F=5\) folds, the approximate \(95\%\) confidence interval is computed using a Student \(t\)-interval:
\[
\bar{m}
\pm
t_{0.975,F-1}
\frac{s_m}{\sqrt{F}}.
\]
Here, \(t_{0.975,F-1}\) is the \(97.5\)th percentile of the \(t\)-distribution with \(F-1\) degrees of freedom. This interval gives a compact indication of fold-level uncertainty for the proposed model. Because it is based on only five folds, it should be interpreted cautiously and used as a robustness summary rather than as a definitive population-level uncertainty estimate \cite{openstax_t_interval}.


\section{Dataset and Fold Summary}
\label{sec:dataset_fold_summary}

The internal experiments were performed on \(N=32\) included patient cases. Each case had one CT volume and one patient-level mucus-score target. The raw target values ranged from \(0.0\) to \(7.5\), with a mean of \(2.688\) and a standard deviation of \(2.313\). The median raw target was \(2.0\), with an interquartile range from \(1.0\) to \(4.0\). The log-transformed target ranged from \(0.000\) to \(0.875\), with a mean of \(0.361\) and a standard deviation of \(0.303\).

The fold sizes were slightly uneven because of the small cohort size. Folds~0 and~1 contained seven patients each, while folds~2, 3, and~4 contained six patients each. Table~\ref{tab:dataset_fold_summary} summarizes the raw target distribution across folds.

\begin{table}[H]
\centering
\caption{Patient counts and raw target distribution across the five internal folds.}
\label{tab:dataset_fold_summary}
\small
\begin{tabular}{cccccc}
\toprule
Fold & Patients & Min & Max & Mean & Median \\
\midrule
0 & 7 & 0.0 & 5.0 & 2.571 & 2.00 \\
1 & 7 & 0.0 & 5.5 & 2.286 & 2.00 \\
2 & 6 & 0.0 & 7.5 & 2.417 & 2.00 \\
3 & 6 & 0.0 & 7.0 & 2.917 & 2.25 \\
4 & 6 & 0.0 & 7.0 & 3.333 & 2.25 \\
\bottomrule
\end{tabular}
\end{table}

The number of slices per CT volume varied from \(47\) to \(716\), with a median of \(241\) slices and an interquartile range from \(151.25\) to \(337\). This variation is one reason why the pipeline uses a slice-selection strategy before model training.


\section{Internal Experiments}
\label{sec:internal_experiments}


\subsection{Baseline Experiments}
\label{sec:baseline_experiments}
The baseline experiments use simpler reference configurations than the main pipeline. They are included to show how the final method compares with more basic slice-handling and input settings. The baseline experiments compare two slice-handling settings using a ResNet18 model with two-window CT input. The uniform-sampling baseline samples slices evenly from the CT volume, while the other ResNet18 model with the same two-window CT input uses the representative slice-selection strategy described in Chapter~4. Because this baseline does not use the full final pipeline configuration, it should be interpreted as a simpler reference setting rather than as a direct ablation of only one component.



\begin{table}[H]
\centering
\caption{Two-window baseline comparison between uniform slice sampling and the representative slice-selection strategy. Values are reported as mean \(\pm\) standard deviation across patient-level folds.}
\label{tab:baseline_results}
\small
\begin{tabular}{lcc}
\toprule
Metric & Uniform sampling & Slice sampling strategy \\
\midrule
MAE & \(1.922 \pm 0.326\) & \(1.760 \pm 0.154\) \\
RMSE & \(2.268 \pm 0.301\) & \(2.223 \pm 0.209\) \\
\(R^2\) & \(-0.009 \pm 0.064\) & \(0.003 \pm 0.109\) \\
Spearman & \(0.202 \pm 0.309\) & \(0.306 \pm 0.330\) \\
\bottomrule
\end{tabular}
\end{table}

The representative slice-selection strategy performed better than uniform sampling on MAE and Spearman correlation, while RMSE and \(R^2\) changed only modestly. This suggests that the selected slices may have provided a more useful input set for patient-level prediction, although the small fold sizes mean that the difference should be interpreted cautiously.


\subsection{Input Representation and Feature Fusion}
\label{sec:input_representation_feature_fusion_results}

Table~\ref{tab:internal_representation_fusion_results} reports the main internal comparison between the two-window input, the three-channel input, and the ORB-BoVW fusion model. These experiments use the main pipeline configuration described in Chapter~4, including the representative slice-selection settings, patient-level aggregation setup, and training configuration summarized in the methodology tables. This makes the rows in this table more directly comparable to each other than to the simpler baseline settings in Table~\ref{tab:baseline_results}.

The models in this comparison were trained using the log-transformed target. For reporting, predictions were transformed back to the raw mucus-score scale so that MAE and RMSE are easier to interpret.

\begin{table}[H]
\centering
\caption{Main internal five-fold results for input representation and feature fusion. Models used the main configuration from Chapter~4, were trained on the log-transformed target, and were transformed back to the raw mucus-score scale for reporting. Values are mean \(\pm\) standard deviation across patient-level folds.}
\label{tab:internal_representation_fusion_results}
\small
\setlength{\tabcolsep}{4pt}
\begin{tabularx}{\textwidth}{@{}>{\raggedright\arraybackslash}Xcccc@{}}
\toprule
Experiment & MAE & RMSE & \(R^2\) & Spearman \\
\midrule
ResNet18, two-window input & \(1.451 \pm 0.395\) & \(1.858 \pm 0.434\) & \(0.126 \pm 0.161\) & \(0.363 \pm 0.239\) \\
ResNet18, three-channel input & \(1.386 \pm 0.456\) & \(1.826 \pm 0.394\) & \(0.186 \pm 0.394\) & \(0.387 \pm 0.256\) \\
\textbf{ResNet18-3CH+BoVW (Proposed model)} & \(\mathbf{1.244 \pm 0.297}\) & \(\mathbf{1.712 \pm 0.505}\) & \(\mathbf{0.236 \pm 0.295}\) & \(\mathbf{0.498 \pm 0.346}\) \\
\bottomrule
\end{tabularx}
\end{table}


The three-channel input performed slightly better than the two-window representation, while \textbf{ResNet18-3CH+BoVW (Proposed model)} gave the strongest overall internal result. This model is treated as the proposed final pipeline because it combines the main three-channel high-pass input with the ORB-BoVW handcrafted feature-fusion branch. The observed pattern is consistent with the high-pass channel and handcrafted local-feature branch adding information that was not fully captured by the CNN representation alone. The interpretation of why this pattern occurs is discussed further in Chapter~6.


\subsubsection{Uncertainty of the Proposed Model}
\label{sec:proposed_model_uncertainty}

To give an additional indication of robustness for the proposed model, Table~\ref{tab:proposed_model_ci} reports approximate 95\% confidence intervals across the five internal validation folds. The intervals are computed from the fold mean and standard deviation using a \(t\)-interval with four degrees of freedom. Since only five folds are available, the intervals should be interpreted cautiously, but they provide a compact summary of the fold-level uncertainty for the final model.

\begin{table}[H]
\centering
\caption{Approximate 95\% confidence intervals for \textbf{ResNet18-3CH+BoVW (Proposed model)} across the five internal validation folds.}
\label{tab:proposed_model_ci}
\small
\begin{tabular}{lcc}
\toprule
Metric & Mean \(\pm\) SD & Approx. 95\% CI \\
\midrule
MAE & \(\mathbf{1.244 \pm 0.297}\) & \([0.875,\ 1.613]\) \\
RMSE & \(\mathbf{1.712 \pm 0.505}\) & \([1.085,\ 2.339]\) \\
\(R^2\) & \(\mathbf{0.236 \pm 0.295}\) & \([-0.130,\ 0.602]\) \\
Spearman & \(\mathbf{0.498 \pm 0.346}\) & \([0.068,\ 0.928]\) \\
\bottomrule
\end{tabular}
\end{table}


\subsection{Backbone Comparison}
\label{sec:backbone_comparison_results}

Table~\ref{tab:backbone_comparison_results} compares ResNet18 and DenseNet121 across the same internal model settings: two-window input, three-channel input, and ORB-BoVW fusion. The ResNet18 rows repeat the corresponding results from Table~\ref{tab:internal_representation_fusion_results} so that the backbone comparison can be read directly. In this comparison, ResNet18 uses the ImageNet-based initialization described in Chapter~4, while DenseNet121 uses pretrained RadImageNet weights. All models were trained on the log-transformed target, and predictions were transformed back to the raw mucus-score scale for reporting.


\begin{table}[H]
\centering
\caption{Backbone comparison across input representation and fusion settings. Values are reported as mean \(\pm\) standard deviation across patient-level folds.}
\label{tab:backbone_comparison_results}
\small
\begin{tabular}{llcc}
\toprule
Setting & Metric & ResNet18 & DenseNet121 \\
\midrule
Two-window input & MAE & \(1.451 \pm 0.395\) & \(1.602 \pm 0.344\) \\
 & RMSE & \(1.858 \pm 0.434\) & \(2.038 \pm 0.483\) \\
 & \(R^2\) & \(0.126 \pm 0.161\) & \(-0.086 \pm 0.168\) \\
 & Spearman & \(0.363 \pm 0.239\) & \(0.208 \pm 0.239\) \\
\midrule
Three-channel input & MAE & \(1.386 \pm 0.456\) & \(1.356 \pm 0.237\) \\
 & RMSE & \(1.826 \pm 0.394\) & \(1.789 \pm 0.384\) \\
 & \(R^2\) & \(0.186 \pm 0.394\) & \(0.084 \pm 0.126\) \\
 & Spearman & \(0.387 \pm 0.256\) & \(0.247 \pm 0.189\) \\
\midrule
ORB-BoVW fusion & MAE & \(1.244 \pm 0.297\) & \(1.422 \pm 0.354\) \\
 & RMSE & \(1.712 \pm 0.505\) & \(1.850 \pm 0.566 \) \\
 & \(R^2\) & \(0.236 \pm 0.295\) & \(0.107 \pm 0.322\) \\
 & Spearman & \(0.498 \pm 0.346\) & \(0.408 \pm 0.358\) \\
\bottomrule
\end{tabular}
\end{table}

DenseNet121 gave slightly lower MAE and RMSE in the three-channel setting, but ResNet18 performed better in the two-window and fusion settings and had stronger rank correlation overall. Therefore, the results do not show a consistent advantage from replacing ResNet18 with DenseNet121.



\section{Transfer Learning Experiments}
\label{sec:transfer_learning_experiments}

The transfer-learning experiments use two source-model variants from the main pipeline. The first variant, \textbf{ResNet18-3CH} corresponds to the ResNet18 slice-based model with the main three-channel high-pass input representation and the main configuration settings described in Chapter~4. The second variant, \textbf{ResNet18-3CH+BoVW}, is the same model but with the ORB-BoVW handcrafted feature-fusion branch added. The fusion variant refers to the \textbf{ResNet18-3CH+BoVW (Proposed model)}, not to a separate standalone handcrafted model.

In the tables, zero-shot means that the source-domain checkpoint is evaluated directly on the target domain without target-domain training. Fine-tuned means that the source-domain checkpoint is used as initialization and then further trained on the target-domain training data before evaluation.

Because the MosMed target differs from the internal mucus-burden target, the transfer tables should be read as evidence about source--target adaptation between CT tasks rather than as clinical external validation for mucus plugging.


\subsection{Mucus-to-MosMed Transfer}
\label{sec:mucus_to_mosmed_transfer}

Table~\ref{tab:transfer_mosmed_target} reports transfer-learning results where the mucus dataset is used as the source domain and MosMed is used as the target domain. In this direction, the target-domain evaluation is performed on the MosMed case-level severity target.

\begin{table}[H]
\centering
\caption{Transfer results from mucus source domain to MosMed target domain.}
\label{tab:transfer_mosmed_target}
\begin{tabular}{lcccc}
\toprule
Transfer setting & MAE & RMSE & \(R^2\) & Spearman \\
\midrule
ResNet18-3CH, zero-shot & 0.903 & 1.070 & -1.62 & -0.027 \\
ResNet18-3CH, fine-tuned on MosMed & 0.460 & 0.643 & 0.04 & 0.320 \\
ResNet18-3CH+BoVW, zero-shot & 0.856 & 1.013 & -0.95 & -0.160 \\
ResNet18-3CH+BoVW, fine-tuned on MosMed & 0.369 & 0.545 & 0.29 & 0.489 \\
\bottomrule
\end{tabular}
\end{table}

The fine-tuned setting performed better than zero-shot evaluation for both model variants. The fine-tuned \textbf{ResNet18-3CH+BoVW (Proposed model)} obtained the lowest MAE and RMSE, and also the highest \(R^2\) and Spearman correlation in this transfer direction.


\subsection{MosMed-to-Mucus Transfer}
\label{sec:mosmed_to_mucus_transfer}

Table~\ref{tab:transfer_mucus_target} reports transfer-learning results where MosMed is used as the source domain and the mucus dataset is used as the target domain. In this direction, the target-domain evaluation is performed on the internal mucus-burden score.

\begin{table}[H]
\centering
\caption{Transfer results from MosMed source domain to mucus target domain.}
\label{tab:transfer_mucus_target}
\begin{tabular}{lcccc}
\toprule
Transfer setting & MAE & RMSE & \(R^2\) & Spearman \\
\midrule
ResNet18-3CH, zero-shot & 1.889 & 2.199 & 0.08 & 0.050 \\
ResNet18-3CH, fine-tuned on mucus & 1.287 & 1.576 & 0.22 & 0.139 \\
ResNet18-3CH+BoVW, zero-shot & 1.802 & 2.097 & 0.10 & 0.261 \\
ResNet18-3CH+BoVW, fine-tuned on mucus & 1.329 & 1.605 & 0.12 & 0.304 \\
\bottomrule
\end{tabular}
\end{table}

The fine-tuned setting performed better than zero-shot transfer on the mucus target for both model variants. The fine-tuned ResNet18-3CH model obtained the lowest MAE, lowest RMSE, and highest \(R^2\), while the fine-tuned \textbf{ResNet18-3CH+BoVW (Proposed model)} obtained the highest Spearman correlation.


\section{Result Summary}
\label{sec:result_summary}
The baseline results show that the representative slice-selection strategy performed better than uniform sampling, with lower MAE and RMSE and a higher Spearman correlation. This gives a simple reference point before comparing the full internal model variants.

In the main internal experiments with the final pipeline configuration, the \textbf{ResNet18-3CH+BoVW (Proposed model)} obtained the lowest MAE and RMSE and the highest \(R^2\) and Spearman correlation among the listed internal experiments. The three-channel input also performed better than the two-window input, although the strongest internal result was obtained when feature fusion was added.

The backbone comparison shows that DenseNet121 did not consistently outperform ResNet18 across the reported settings. DenseNet121 gave slightly lower MAE and RMSE in the three-channel setting, but ResNet18 performed better in the two-window and fusion settings and had stronger rank correlation overall.

In the transfer-learning experiments, the fine-tuned settings performed better than zero-shot transfer in both directions. The transfer results are reported as average performance across the five folds. In the mucus-to-MosMed direction, the fine-tuned \textbf{ResNet18-3CH+BoVW (Proposed model)} performed best across all reported metrics for the MosMed severity target. In the MosMed-to-mucus direction, the fine-tuned ResNet18-3CH model obtained the lowest MAE, lowest RMSE, and highest \(R^2\), while the fine-tuned \textbf{ResNet18-3CH+BoVW (Proposed model)} obtained the highest Spearman correlation. These findings describe transfer behaviour between different CT prediction tasks and should not be interpreted as external mucus-specific validation.

Several \(R^2\) values are close to zero or negative, especially in the baseline and zero-shot transfer settings. These values are kept in the tables because they show the difficulty of the task and the limited usefulness of direct transfer without target-domain adaptation. They also support interpreting the results as exploratory patient-level prediction results rather than definitive clinical evidence.

The deeper interpretation of these results, including the role of slice selection, high-pass input representation, feature fusion, backbone choice, and transfer direction, is discussed in Chapter~6.

\chapter{Discussion}


\section{Overview of Main Findings}
The main goal of this thesis was to investigate whether a slice-based CT pipeline could predict patient-level mucus burden from volumetric chest CT scans under weak supervision. The results show that this is possible to some extent, but also that the task remains difficult. The best internal model did not achieve perfect prediction, but the pattern across the experiments suggests that several parts of the pipeline may have contributed useful information. This section summarizes the findings in relation to RQ1--RQ5 before the following sections discuss each component in more detail.

The baseline experiments showed that the representative slice selection performed better than uniform slice sampling. This is an important result because the model does not process the full CT volume directly. It depends on which slices are selected before feature extraction. The comparison suggests that selecting slices based on central restriction, informativeness filtering, and diversity sampling may give the model a better input set than simply sampling slices evenly through the volume.

The main internal experiments also showed better results for the three-channel high-pass representation than for the two-window input, and the strongest internal result was obtained when ORB-BoVW feature fusion was added. The \textbf{ResNet18-3CH+BoVW (Proposed model)} achieved the lowest MAE and RMSE, and the highest \(R^2\) and Spearman correlation among the reported internal configurations. This suggests that the learned CNN representation and the handcrafted local-feature branch may have captured some complementary information.

The backbone comparison did not show that a larger or more domain-specific pretrained backbone was automatically better. DenseNet121 with RadImageNet initialization gave slightly better MAE and RMSE in the three-channel setting, but ResNet18 performed better in the two-window and fusion settings and had stronger rank correlation overall. In this dataset ResNet18 appeared to be the more reliable backbone across different configurations.

The transfer-learning experiments showed that fine-tuning was more useful than zero-shot transfer in both transfer directions. This is expected because the internal mucus dataset and MosMed differ in both dataset properties and target definition. The transfer results therefore support the idea that pretrained or source-trained representations can be useful under domain shift, but they still need target-domain adaptation to perform well.

Overall, the results support the main design direction of the thesis. In the tested configurations, patient-level CT prediction was associated with better performance when slice selection, richer input representation, and feature fusion were included. At the same time, the findings remain exploratory because the internal cohort is small, supervision is only available at patient level, and the transfer setup uses a different external target. These limitations are discussed later, but they are important context for reading the component-wise patterns.

These findings answer the research questions from Chapter~1 in a cautious way. For RQ1, the results show that patient-level mucus-burden prediction from CT is possible to some extent under weak supervision, but still difficult. RQ2 is addressed by the baseline comparison, where the representative slice selection performed better than uniform sampling. RQ3 is addressed by the input-representation and feature-fusion results, where the high-pass and ORB-BoVW configurations were associated with better internal model results. RQ4 is addressed by the backbone comparison, where ResNet18 was more reliable overall than DenseNet121 in this setup. Finally, RQ5 is addressed by the transfer-learning experiments, where fine-tuning was more useful than zero-shot transfer across the two source--target directions. The following sections discuss these points in more detail by linking the results back to the main pipeline components and their limitations.


\section{Ablation-Style Interpretation of Pipeline Components}
The results can be read as an ablation-style progression from simpler settings toward the final pipeline. This section mainly supports RQ2, RQ3, and RQ4 by examining what happens when important components are added or compared, first the slice-selection strategy, then the high-pass input channel, then handcrafted feature fusion, and finally the choice of CNN backbone. These comparisons help show which design choices appeared most useful in the final patient-level prediction setup.

At the same time, the comparisons should not be interpreted as perfect single-factor ablations in every case. Some comparisons in this thesis are ablation-style rather than strict single-factor ablations, because multiple settings may differ between configurations. The results therefore indicate which configurations were associated with better performance, but they do not always isolate the causal effect of each component. This is especially important because the dataset is small and the validation folds contain only a limited number of patients.


\subsection{Representative Slice Selection Strategy}
The first component to consider is the representative slice-selection strategy, which directly relates to RQ2. Since the model does not process the full CT volume directly, the selected slices define the image evidence that the CNN receives. If this step selects many uninformative slices, the later model stages have less useful input, even if the CNN architecture itself is unchanged.

The baseline comparison showed that representative slice selection performed better than uniform sampling. With uniform sampling, the ResNet18 two-window baseline obtained an MAE of \(1.922\), RMSE of \(2.268\), \(R^2\) of \(-0.009\), and Spearman correlation of \(0.202\). When the representative slice selection was used, the corresponding values were an MAE of \(1.760\), RMSE of \(2.223\), \(R^2\) of \(0.003\), and Spearman correlation of \(0.306\). The difference is not dramatic for every metric, but the direction is consistent.

This is consistent with the slice-selection strategy giving the model a better patient-level slice set than simply sampling evenly through the scan. This makes sense for the task, because mucus-related evidence is not expected to be equally visible in every axial slice. Some slices may contain little relevant lung content, while others may contain more useful local information. The central restriction, informativeness filtering, and diversity sampling may therefore reduce the chance that the model spends much of its capacity on irrelevant or weak slices.

Another advantage of the representative strategy is that it reduces dependence on which axial positions happen to be sampled. A uniform strategy may include useful slices, but it can also miss informative regions or include many slices with little relevant lung content, especially when the number of selected slices is limited. This can make the baseline more sensitive to scan length, slice spacing, and the exact sampled positions. In contrast, the representative strategy uses fixed rules for central restriction, informativeness filtering, and diversity sampling, so the selected set is guided by image content rather than by axial position alone.

However, this should be read as patient-level slice selection, not localization. The strategy is heuristic, unsupervised with respect to mucus burden, and does not use slice-level mucus labels. Its role is more modest: it creates a more useful and stable subset of slices for patient-level prediction. Within that role, the baseline comparison supports keeping the representative slice-selection stage in the final pipeline.


\subsection{Input Representation and High-Pass Channel}
The next step in the pipeline was the move from the two-window input to the three-channel input, which is part of RQ3. The two-window representation already gives the model two useful views of the CT slice, one focused on lung-window information and one focused on mediastinal-window information. The three-channel version keeps these two channels and adds a high-pass channel. The purpose of the high-pass channel is to emphasize local intensity changes and fine image structure.

The internal results indicate that the three-channel input outperformed the two-window input. The ResNet18 two-window model achieved a mean absolute error (MAE) of \(1.451\), root mean square error (RMSE) of \(1.858\), coefficient of determination (\(R^2\)) of \(0.126\), and Spearman correlation of \(0.363\). With the high-pass third channel, the ResNet18 three-channel model used the same main pipeline configuration and achieved an MAE of \(1.386\), RMSE of \(1.826\), \(R^2\) of \(0.186\), and Spearman correlation of \(0.387\). Although the improvement is smaller than the difference observed with \textbf{ResNet18-3CH+BoVW (Proposed model)}, it remains consistent across all reported metrics.

This suggests that the high-pass channel may have added useful information beyond the original two CT windows. A possible explanation is that the high-pass channel makes local changes in the image more visible to the CNN. This can be useful in a task where the relevant evidence may be small, local, and not evenly distributed across the lungs. The implementation of lung-mask suppression also helps focus this third-channel response inside the approximate lung region instead of letting high-frequency structures outside the lungs dominate the input.

At the same time, the high-pass channel should not be interpreted as a direct mucus detector. It highlights local image structure rather than clinically verified mucus findings. The safer interpretation is that it gave the model an additional view of the selected slices, and that this extra view was useful for patient-level prediction in the reported experiments.


\subsection{Handcrafted ORB-BoVW Feature Fusion}
The strongest internal result for RQ3 was obtained when the CNN features were combined with the ORB-BoVW handcrafted feature branch. Compared with the three-channel CNN-only model, the \textbf{ResNet18-3CH+BoVW (Proposed model)} reduced the MAE from \(1.386\) to \(\mathbf{1.244}\), reduced the RMSE from \(1.826\) to \(\mathbf{1.712}\), and increased the Spearman correlation from \(0.387\) to \(\mathbf{0.498}\). This was the clearest performance gain associated with a configuration change in the internal experiments, and it suggests that the handcrafted branch may have added information that the CNN did not fully capture on its own.

One possible reason is that ORB-BoVW describes the image in a different way than the CNN. The CNN learns features directly from the training data, while the ORB-BoVW branch summarizes local keypoint patterns using a fixed visual vocabulary. In a larger dataset the CNN might learn many of these local structures by itself. In this thesis the internal cohort is small, so the handcrafted branch may have helped by adding a more stable description of local image texture and structure. This fits the general idea behind the fusion setup, the handcrafted features are not meant to replace the CNN, but to give the model another type of signal.

The result is also interesting because mucus-burden prediction is likely influenced by local and unevenly distributed findings. A patient-level score can depend on small regions in a subset of slices rather than on the whole scan equally. The high-pass channel already gives the CNN an input that emphasizes local intensity changes, but the ORB-BoVW branch adds another summary of local image patterns. The fusion result suggests that these two views were complementary in this experimental setting.

At the same time, ORB-BoVW keypoints are not mucus-specific. The branch describes local image structure through keypoints and visual words, so its benefit should be interpreted as additional patient-level signal rather than direct mucus detection.

There is also a trade-off. Adding ORB-BoVW fusion increases the pipeline's complexity. It requires fitting a visual vocabulary within each training fold, caching or computing handcrafted features, and making sure validation data are not used when building the vocabulary. This increases the number of moving parts compared with a CNN-only model. Because the fusion configuration had better values for all reported internal metrics, the extra complexity appears justified in this thesis setup.

Overall, the ORB-BoVW result supports the idea that handcrafted features can still be useful in small medical-imaging datasets. The best interpretation is not that handcrafted features are better than learned features, but that they can complement learned CNN features when the amount of training data is limited and the target is only available at patient level.


\subsection{Backbone Comparison}
The backbone comparison addresses RQ4 by comparing whether a different CNN feature extractor was associated with better patient-level prediction. In this thesis, ResNet18 and DenseNet121 were compared across the same main settings, the two-window input, the three-channel input, and the ORB-BoVW fusion. DenseNet121 is a deeper architecture and was initialized with RadImageNet weights, so it was reasonable to expect that it might provide stronger medical-image features. However, the results did not show a clear overall advantage from switching to DenseNet121.

For the two-window input, ResNet18 performed better than DenseNet121 across all reported metrics. ResNet18 achieved an MAE of \(1.451\), RMSE of \(1.858\), and Spearman correlation of \(0.363\), while DenseNet121 obtained an MAE of \(1.602\), RMSE of \(2.038\), and Spearman correlation of \(0.208\). This suggests that the larger backbone was not beneficial in the simpler two-window setting.

The three-channel setting was more mixed. DenseNet121 gave slightly lower MAE and RMSE than ResNet18, with MAE decreasing from \(1.386\) to \(1.356\) and RMSE from \(1.826\) to \(1.789\). However, ResNet18 still had better \(R^2\) and Spearman correlation. This means that DenseNet121 reduced the average error slightly in this setting, but it did not rank patients as well as ResNet18. Since Spearman correlation is important for understanding whether the model orders patients correctly by predicted burden, this difference matters.

The fusion setting again favoured ResNet18. \textbf{ResNet18-3CH+BoVW (Proposed model)} achieved the best overall internal result, while DenseNet121 performed worse on MAE, RMSE, \(R^2\), and Spearman. If DenseNet121 had been clearly better as a backbone in this setup, one would expect it to perform well in this setting as well, but that was not observed.

One possible explanation is that DenseNet121 may have been harder to adapt in this small dataset. The internal cohort contains only 32 patient cases, and each validation fold contains only a few patients. In such a setting, a larger or more complex backbone does not automatically lead to better generalization. ResNet18 is simpler, and this may have made it more stable when combined with the selected slice representation, aggregation strategy, and handcrafted fusion branch.

This does not mean that ResNet18 is generally better than DenseNet121 for CT analysis. It only means that, in this experimental setup, DenseNet121 did not provide a consistent advantage. The results suggest that backbone choice should be evaluated together with the rest of the pipeline, rather than assuming that a deeper or radiology-pretrained model will always perform better. For this thesis, ResNet18 was therefore kept as the main backbone because it gave the strongest and most consistent results across the reported configurations.


\subsection{Aggregation and Patient-Level Prediction}
Aggregation was an important part of the pipeline because the model produces slice-level scores, while the target is defined at patient level. During development, this was approached step by step. The simplest option was mean aggregation, where all selected slice scores contribute equally to the final patient prediction. This is stable and easy to interpret, but it assumes that all selected slices are equally informative.

The next option was top-\(k\) aggregation. This was motivated by the idea that mucus-related evidence may only be visible in a subset of slices. Instead of averaging all slice scores, top-\(k\) aggregation focuses on the highest predicted slice scores. This can be useful when the patient-level target is driven by the most abnormal or informative slices, but it also risks ignoring broader information from the rest of the scan.

The final hybrid aggregation strategy was built from these two ideas. It combines the mean prediction with the top-\(k\) prediction, so the final patient-level estimate keeps information from all selected slices while still giving extra weight to the slices with stronger predicted evidence. This made sense for the task because mucus burden is patient-level, but the visual evidence can be local and unevenly distributed.

This aggregation step also reflects the weakly supervised setup. The slice-level scores are useful only after being combined into a patient-level prediction; they are not independently verified slice findings. The purpose of aggregation is therefore to connect these intermediate scores to the available patient-level target.

Although the final pipeline uses hybrid aggregation, this thesis does not present a separate result table comparing mean, top-\(k\), and hybrid aggregation. Therefore, the discussion should not claim that hybrid aggregation was proven to be optimal. A more careful interpretation is that hybrid aggregation was chosen because it combines the stability of mean aggregation with the more focused behaviour of top-\(k\) aggregation. A full aggregation ablation would be a useful direction for future work.


\subsection{Target Transformation and Reporting Scale}
The target scale also affects how the model is trained and how the results should be interpreted. In the baseline experiments, the raw mucus-burden score was used directly. In the main internal experiments and transfer-learning experiments, the models were trained using the log-transformed target, and predictions were transformed back to the raw mucus-score scale for reporting.

This choice was made because the raw mucus-burden scores are not evenly distributed. Most patients have relatively low or moderate values, while a few patients have higher scores. In this situation, training directly on the raw scale can make the model more sensitive to larger target values. The log transformation can reduce the influence of these larger values and may make the regression problem easier to optimize.

At the same time, reporting MAE and RMSE on the log scale would be harder to interpret clinically and practically. For that reason, predictions were transformed back to the raw score scale before reporting the main error metrics. This makes the reported MAE and RMSE easier to understand, because they are expressed in the same units as the original mucus-burden score.

The log transformation should also be interpreted carefully in the ablation argument. The raw-target baseline and the log-target main experiments do not differ only in target scale, they also differ in other parts of the configuration. Therefore, the results do not isolate the effect of the log transformation alone. It is more accurate to say that the log target was part of the final training setup that gave the best internal results, rather than claiming that the log transformation alone caused the improvement.

The target transformation was part of a training setup intended to handle a small and unevenly distributed regression target, while the back-transformation kept the reported results understandable on the original score scale. This balance is important because the model should be optimized in a stable way, but the results still need to be interpretable in terms of patient-level mucus burden.


\section{Transfer Learning and Domain Shift}
The transfer-learning experiments address RQ5 by asking whether external CT data with a different target can support representation transfer in a small-cohort setting. This matters because the internal mucus dataset is small, so it is natural to ask whether a model trained on another CT task can provide a useful starting point.

The transfer experiments used two model variants from the main pipeline. The first was ResNet18-3CH, which refers to the ResNet18 model with the main three-channel high-pass input. The second was \textbf{ResNet18-3CH+BoVW (Proposed model)}, which adds the ORB-BoVW handcrafted feature-fusion. They include the main configuration for the slice-based pipeline, input representation, and patient-level aggregation described in Chapter~4.

The main pattern is clear, fine-tuning was more useful than zero-shot transfer in both directions. This suggests that the models could reuse some information from the source domain, but they still needed target-domain adaptation. This is expected because the internal mucus dataset and MosMed are both based on chest CT, but they do not have the same target. The internal task predicts mucus burden, while MosMed uses a case-level CT severity label related to COVID-19 lung involvement.


\subsection{Zero-Shot Transfer}
Zero-shot transfer means that a model trained on the source domain is evaluated directly on the target domain without any target-domain training. In this thesis, zero-shot transfer was limited in both directions.

In the mucus-to-MosMed direction, the ResNet18-3CH zero-shot model obtained an MAE of \(0.903\), RMSE of \(1.070\), \(R^2\) of \(-1.62\), and Spearman correlation of \(-0.027\). The \textbf{ResNet18-3CH+BoVW (Proposed model)} zero-shot model was slightly better on MAE and RMSE with values of \(\mathbf{0.856}\) and \(\mathbf{1.013}\), but it still had a negative \(R^2\) of \(\mathbf{-0.95}\) and a negative Spearman correlation of \(\mathbf{-0.160}\). This shows that even the proposed model trained on the mucus task did not transfer well to the MosMed severity task without adaptation.


The MosMed-to-mucus direction showed a similar pattern. The ResNet18-3CH zero-shot model obtained an MAE of \(1.889\), RMSE of \(2.199\), \(R^2\) of \(0.08\), and Spearman correlation of \(0.050\). The \textbf{ResNet18-3CH+BoVW (Proposed model)} zero-shot model performed better, with an MAE of \(\mathbf{1.802}\), RMSE of \(\mathbf{2.097}\), \(R^2\) of \(\mathbf{0.10}\), and Spearman correlation of \(\mathbf{0.261}\). This suggests that the BoVW fusion branch transferred somewhat better than the non-fusion variant in this direction, but the zero-shot setting was still weaker than fine-tuning.

These results are not surprising. Zero-shot transfer assumes that the source-domain model can be applied directly to a different dataset and a different target definition. Even though both datasets contain chest CT scans, the model is not solving the same prediction task in both cases. The weak zero-shot results therefore support the idea that representation reuse alone is not enough for this problem.


\subsection{Fine-Tuning}
The fine-tuned setting performed better than zero-shot transfer in both directions. This means that the source-domain model was more useful as an initialization than as a final model applied directly to the target domain.

In the mucus-to-MosMed direction, the fine-tuned ResNet18-3CH model had a lower MAE than the zero-shot setting, from \(0.903\) to \(0.460\), and a lower RMSE, from \(1.070\) to \(0.643\). The \textbf{ResNet18-3CH+BoVW (Proposed model)} showed the same pattern, with MAE changing from \(0.856\) to \(\mathbf{0.369}\), and RMSE from \(1.013\) to \(\mathbf{0.545}\). The fine-tuned \textbf{ResNet18-3CH+BoVW (Proposed model)} also achieved the strongest MosMed-target result, with \(R^2=\mathbf{0.29}\) and Spearman correlation \(\mathbf{0.489}\).


In the MosMed-to-mucus direction, the fine-tuned settings also had better error metrics than the zero-shot settings. The ResNet18-3CH model changed from an MAE of \(1.889\) to \(1.287\), and from an RMSE of \(2.199\) to \(1.576\). The \textbf{ResNet18-3CH+BoVW (Proposed model)} changed from an MAE of \(1.802\) to \(\mathbf{1.329}\), and from an RMSE of \(2.097\) to \(\mathbf{1.605}\). In this direction, the fine-tuned ResNet18-3CH model gave the lowest MAE, lowest RMSE, and highest \(R^2\), while the fine-tuned \textbf{ResNet18-3CH+BoVW (Proposed model)} gave the highest Spearman correlation.


The main interpretation is that source-domain training can provide a useful starting point, but the model still needs to adapt to the target dataset. This is especially relevant in small medical-imaging datasets, where training from scratch is difficult, but direct zero-shot transfer is usually too weak. Fine-tuning gives a middle ground, the model starts from features learned on another CT task, and then the later training adjusts those features to the target-domain label.


\subsection{Transfer Direction and Dataset Mismatch}
The transfer results should be interpreted with the dataset mismatch in mind. MosMed is not an external mucus-burden validation dataset. Its target is a different case-level CT severity label, while the internal target is a mucus-burden score. This means that the transfer experiments test whether CT representations can be adapted between related tasks, and not whether the mucus model generalizes directly to an external mucus-scored cohort.

This target mismatch helps explain why fine-tuning was needed. A model trained on mucus burden may learn useful CT features, but those features still need to be connected to the MosMed severity label. The same is true in the opposite direction, a MosMed-trained model may learn features related to lung involvement, but it must be adapted before it can predict mucus burden.

The direction of transfer also matters. In the mucus-to-MosMed direction, the fine-tuned settings performed better than the zero-shot settings for both ResNet18-3CH and \textbf{ResNet18-3CH+BoVW (Proposed model)}. In this direction, the proposed model was strongest across all reported metrics. In the MosMed-to-mucus direction, the fine-tuned settings also performed better than the zero-shot settings, but the best model depended on the metric. ResNet18-3CH had the lowest MAE, lowest RMSE, and highest \(R^2\), while \textbf{ResNet18-3CH+BoVW (Proposed model)} had the highest Spearman correlation. This suggests that transfer behaviour is not fully symmetric between the two datasets.

Overall, the transfer experiments answer the transfer-related research question in a cautious way. They show that external CT data and source-domain training can be useful for representation transfer, but mainly when followed by target-domain fine-tuning. They also show that zero-shot transfer is not enough when the dataset and target definition change. For this thesis, transfer learning should therefore be viewed as a useful adaptation strategy, not as external clinical validation of mucus-burden prediction.



\section{Weak Supervision and Interpretation of Slice-Level Evidence}
A central challenge in this thesis is that the available labels are defined at patient level, while the model is trained on selected CT slices. This point is central to RQ1 and RQ2, because the thesis evaluates patient-level prediction while also asking how the selected slice set affects that prediction. During training, each selected slice from patient \(i\) is paired with the same patient-level target \(y_i\). For example, if a patient has a mucus-burden score of \(5\), then all selected slices from that patient are trained with target value \(5\), even though the actual mucus evidence may only be visible in some of those slices.

This is the weakly supervised part of the setup. The label is correct for the patient, but it is not a verified label for each individual slice. A slice-level score \(s_{i,j}\) may be useful for the final patient-level prediction, but it should not be treated as a true slice-level mucus measurement.

At validation and testing time, the model produces slice-level scores for the selected slices. These scores are aggregated into one patient-level prediction \(\hat{y}_i\), which is compared with the true patient-level target \(y_i\). The evaluation is therefore patient-level, not slice-level. The pipeline can be judged by how well it predicts the patient score, but not by whether individual slice scores correctly localize mucus.

This also affects how the slice-selection strategy and the high-pass input representation should be understood. These components may help the model build a better patient-level prediction, but they should not be read as clinical explanation maps. Different slice-score patterns could also lead to similar patient-level predictions; a patient score could be driven by a few high-scoring slices or by many moderately scoring slices. Without slice-level annotations, it is not possible to know which pattern is more clinically correct.

For this reason, the results should be interpreted as patient-level prediction results, not localization results. The pipeline suggests that selected slices and local image features contain useful information for predicting mucus burden, but it does not prove that the model has learned to identify mucus plugs directly. A future version of the project would need slice-level or region-level annotations to evaluate that question properly.


\section{Reliability, Generalization, and Limitations}
The main limitation of this thesis is the small size of the internal dataset. The internal experiments were performed on 32 patient cases, with only six or seven patients in each validation fold. This makes the results sensitive to individual cases. A single difficult or unusual patient can have a large effect on fold-level metrics, especially for RMSE and \(R^2\). The standard deviations reported in Chapter~5 should therefore be taken seriously, not treated as small details.

The low and sometimes negative \(R^2\) values are also important for interpreting the results honestly. They show that the models do not consistently explain a large share of target variance. Some settings perform close to, or worse than, a simple mean-prediction reference. This is not unexpected given the small internal cohort, weak patient-level supervision, possible noise in the mucus-score target, and the difficulty of mapping selected CT slices to one patient-level burden score. MAE and RMSE are therefore important complementary metrics because they show the size of the prediction error on the target scale, while Spearman correlation helps indicate whether patients are ranked in a meaningful order.

The small cohort also limits how strongly the ablation-style findings can be generalized. The results suggest that representative slice selection strategy, high-pass input representation, and ORB-BoVW fusion were useful in this dataset. However, the same performance patterns may not appear in the same way on a larger or more diverse cohort. This is especially relevant for the handcrafted fusion branch. It performed well here, but it also adds more preprocessing and feature-engineering steps, which may behave differently on data from other scanners, hospitals, or patient groups.

The patient-level supervision also limits interpretability. Even when a model predicts the patient score reasonably, it cannot be checked directly against slice-level mucus findings in this dataset. This is the main reason the thesis avoids localization claims.

The transfer-learning experiments add a second generalization limit. MosMed is useful for studying transfer between related CT tasks, but its target is not mucus burden. These results therefore describe domain-shift behaviour, not external mucus-specific validation. A true external validation would require another dataset with the same or a closely comparable mucus-burden target.

The negative \(R^2\) values in the zero-shot transfer experiments are especially relevant here. They indicate that direct source-to-target transfer can be worse than predicting the target-domain mean when the dataset and label definition change. This reinforces the main transfer-learning interpretation that source-trained CT representations may be useful after fine-tuning, but direct zero-shot reuse is not reliable in this setting.

There are also limitations in the experimental coverage. The thesis compares several important components, but it does not provide a complete ablation of every hyperparameter. For example, the final aggregation strategy was motivated by development experiments with mean and top-\(k\) aggregation, but Chapter~5 does not report a separate aggregation comparison table. Similarly, several third-channel variants were explored during development, including high-pass, DoG, LoG, and DCT-based representations. The high-pass representation was kept as the main three-channel setup because it gave the strongest and most stable results in these development runs. However, these exploratory comparisons were not reported as full formal result tables in Chapter~5, since the results chapter focuses on the main pipeline configurations. This means that the thesis supports the high-pass choice as a development-driven design decision, but does not present a complete formal ablation of every third-channel variant. 

The pipeline is based on two-dimensional slice processing rather than full three-dimensional volume modelling. This made the task more manageable with the available data and computational resources, but it also means that the model may miss some spatial context across neighbouring slices. The representative slice-selection strategy and aggregation stages partly address this by selecting multiple slices and combining their predictions, but they are not the same as modelling the full 3D structure of the CT volume.

Overall, the results are promising, but they should be viewed as preliminary. The thesis shows that the proposed pipeline can learn useful patient-level signals from CT under weak supervision, and that several design choices were associated with better performance within the available data. At the same time, stronger conclusions would require a larger cohort, external mucus-specific validation, and more detailed annotation for evaluating slice-level or regional evidence.


\section{Future Work}
Several directions could make the pipeline stronger in future work. The most important one is to evaluate the method on a larger internal cohort. The current dataset contains only 32 patient cases, which limits how confidently the results can be generalized. A larger dataset would make the cross-validation results more stable and would also make it easier to test whether the performance patterns associated with slice selection, high-pass input, and ORB-BoVW fusion still hold when more variation is present. 

A second important direction is external validation on another dataset with mucus-burden labels. Since MosMed was used for transfer and domain-shift analysis, it cannot serve this role. A true external validation would require an independent cohort where mucus burden is scored in a similar way to the internal dataset.

More detailed annotations would also be valuable. Slice-level or region-level mucus annotations would make it possible to test whether high-scoring slices correspond to mucus-related areas, rather than only supporting the final patient-level score.

The aggregation strategy could also be extended in future work. The final pipeline used a hybrid aggregation rule based on mean and top-\(k\) aggregation, which was chosen because it gave a good balance between using information from all selected slices and emphasizing the most informative slice predictions. Future work could present a more detailed comparison of aggregation choices, including mean, max, top-\(k\), hybrid aggregation, attention-based aggregation, and multiple-instance learning approaches. This would make the aggregation part of the pipeline more transparent and help show when each strategy is most useful.

The input representation could also be studied further. The high-pass channel was selected as the main third-channel configuration because it performed best during development and gave a simple way to emphasize local image structure. Other variants, such as DoG, LoG, and DCT-based filtering, were explored but not reported as full result tables in Chapter~5. Future work could present these comparisons more formally, together with the role of lung-mask suppression, the choice of high-pass source channel, and the high-pass scale. This would make the input-representation design more complete and easier to compare.

The ORB-BoVW fusion was associated with the strongest main internal results, but it adds extra complexity to the pipeline. Future work could test different vocabulary sizes, keypoint sources, descriptor types, or simpler handcrafted alternatives. It would also be useful to investigate whether the fusion branch remains helpful when the dataset becomes larger, since a CNN may be able to learn more local structure by itself when more training data are available.

Finally, future work could explore models that use more spatial context. The current pipeline is based on selected two-dimensional slices, which makes the problem manageable with limited data and computation. However, mucus burden is a volumetric finding, and neighbouring slices may contain useful context. A future model could therefore use 2.5D inputs, short slice sequences, recurrent aggregation, attention-based slice pooling, or full 3D CNNs if enough data are available. These approaches could make greater use of volumetric information while still keeping the patient-level prediction task central.


\section{Concluding Discussion}
This thesis investigated a weakly supervised, slice-based pipeline for predicting patient-level mucus burden from chest CT. The results show that the task is possible to approach with the available data, but also that it remains challenging. The best internal model did not produce perfect predictions, but the experimental progression showed that several richer configurations performed better than simpler settings.

The most important internal finding was that the \textbf{ResNet18-3CH+BoVW (Proposed model)} performed best among the reported configurations. The representative slice selection performed better than uniform sampling, the high-pass three-channel configuration performed better than the two-window input, and adding ORB-BoVW feature fusion gave the strongest internal result. These findings are consistent with the idea that the problem depends not only on the CNN backbone itself, but also on how the CT volume is sampled, represented, and combined with complementary local features.

The transfer-learning experiments also gave a useful result. Zero-shot transfer was limited, while the fine-tuned settings performed better in both transfer directions. This suggests that source-domain CT representations can be useful, but they need to be adapted to the target dataset and target label. This is adaptation evidence, not external mucus validation.

A key point throughout the thesis is that the work is weakly supervised. The slice-level scores are intermediate outputs trained from patient-level labels, so the pipeline should be understood as patient-level prediction rather than mucus localization.

Overall, the thesis shows that a carefully designed slice-based CT pipeline can learn useful patient-level signals for mucus-burden prediction in a small-cohort setting. The results support the use of representative slice selection, high-pass input representation, and handcrafted feature fusion as meaningful parts of the pipeline. However, the findings are still preliminary because of the small dataset, weak supervision, and lack of external mucus-specific validation. The method is therefore best viewed as a promising experimental framework that can be strengthened through larger datasets, more detailed annotations, and broader validation.

\chapter{Conclusion}


\section{Thesis Summary}
This thesis investigated automated prediction of patient-level mucus plug burden from lung CT scans. The task was formulated as weakly supervised regression, each CT examination had one patient-level mucus-burden score, while the model operated on selected axial slices without slice-level mucus annotations. This made the problem different from direct plug detection or segmentation. The goal was not to build a clinically ready system, but to test whether a constrained slice-based pipeline could learn useful patient-level signals from the available data.

The implemented pipeline combined stored CT volume loading, intensity handling, representative slice selection, CT input representation, CNN feature extraction, optional ORB-BoVW feature fusion, slice-level prediction, and patient-level aggregation. The main internal evaluation used five-fold cross-validation on the internal mucus-score cohort, while the transfer experiments used MosMed as a separate CT dataset for studying representation transfer and domain shift. Because MosMed has a COVID-related severity target rather than a mucus-burden target, it was not treated as external validation of mucus prediction.

The results answer the research questions in a cautious but useful way. For RQ1, the internal experiments showed that patient-level prediction is possible to approach with the proposed weakly supervised pipeline, but the low and sometimes negative \(R^2\) values also show that the task remains difficult. For RQ2, representative slice selection performed better than uniform slice sampling in the baseline comparison, suggesting that the way slices are selected matters when the full CT volume is not processed end-to-end. For RQ3, the three-channel high-pass input and ORB-BoVW fusion were associated with the strongest internal results, with the ResNet18 three-channel fusion model performing best among the reported internal configurations. For RQ4, DenseNet121 did not consistently outperform ResNet18, even though it gave slightly better error metrics in one three-channel setting. For RQ5, the transfer experiments showed that fine-tuning was more useful than zero-shot transfer, while also showing the limits of direct transfer when the target definition and domain change.

These findings support the main idea of the thesis, that predicting mucus burden from CT scans at the patient level can be explored using a slice-based weakly supervised pipeline, and that some design choices seem helpful in this experimental setup. At the same time, the results should still be seen as exploratory. The cohort is small, the labels are only available at the patient level, and the transfer experiments are not a substitute for validation on an independent dataset with mucus scores.


\section{Main Contributions}
The first contribution is the formulation of mucus plug burden prediction as a patient-level weakly supervised regression problem. This framing is important because it keeps the prediction target aligned with the available clinical score while acknowledging that the image evidence is distributed across a volumetric CT examination. It also makes clear that slice-level outputs from the model are intermediate signals, not verified mucus-localization results.

The second contribution is the implementation of a complete and reproducible slice-based CT pipeline. The pipeline separates the main stages of the problem: preprocessing, representative slice selection strategy, input representation, feature extraction, aggregation, and model selection. This separation made it possible to test the role of different parts instead of treating the model as one whole system.

The third contribution of this thesis is comparing different design choices when only a small amount of data is available. For example, representative slice selection was compared with uniform sampling, two-window input was compared with a three-channel high-pass method, and ResNet18 was compared with DenseNet121. The study also tested ORB-BoVW fusion as a handcrafted feature branch. These comparisons are not perfect ablation studies, but they still help show which choices gave better results in this thesis.

The fourth contribution is the transfer-learning analysis using MosMed. These experiments did not validate mucus prediction externally, but they did give useful information about source--target adaptation between CT tasks. The difference between zero-shot transfer and fine-tuning showed that source-trained CT representations can be helpful, but that direct reuse across a different target definition is limited.


\section{Future Work}
The most important next step is evaluation on a larger mucus-scored cohort. With only 32 internal cases, each validation fold is small, and individual patients can strongly affect the reported metrics. A larger cohort would make the results more stable and would make it easier to test whether the observed patterns for slice selection, high-pass input, feature fusion, and backbone choice remain consistent.

External validation should also be performed on an independent dataset with a comparable mucus-burden target. The MosMed experiments were useful for studying transfer and domain shift, but they cannot answer whether the model generalizes to a separate mucus-scored population. A true external validation dataset would be needed before making stronger claims about generalization.

Future work should also include more detailed annotations if they become available. Slice-level or region-level mucus labels would make it possible to evaluate whether high-scoring slices correspond to actual mucus findings. This would strengthen interpretability and would help separate patient-level prediction from localization, which this thesis could not do.

Several methodological directions are also worth extending. A more controlled aggregation study could compare mean, top-\(k\), hybrid, and attention-based pooling in a cleaner way. The input-representation variants could be evaluated more systematically, including different high-pass settings and other frequency-based channels. The ORB-BoVW branch should also be tested in larger data to see whether handcrafted local features remain useful when the CNN has more examples to learn from. Finally, models that use more spatial context, such as 2.5D inputs, short slice sequences, or volumetric models, may be useful if enough data are available to support them.


\section{Concluding Remarks}
This thesis shows that automated patient-level mucus-burden prediction from CT is a meaningful but difficult machine-learning problem. The best internal results came from a pipeline that combined representative slice selection strategy, a three-channel high-pass CT representation, ResNet18 feature extraction, ORB-BoVW fusion, and patient-level aggregation. These findings suggest that careful input construction and complementary feature representations matter in a small weakly supervised CT setting.

The work also shows why the problem should be treated cautiously. Low and negative \(R^2\) values in several settings, the small internal cohort, patient-level-only labels, and the absence of external mucus-specific validation all limit the strength of the conclusions. The transfer experiments further show that CT representations do not transfer reliably without adaptation when the dataset and target definition change.

The final takeaway is therefore balanced. The proposed pipeline is not a clinical mucus-plug detection system, and it does not prove that the model localizes mucus plugs. What it does provide is a reproducible experimental framework and evidence that patient-level mucus-burden prediction from CT can benefit from thoughtful slice selection, input representation, feature fusion, and fine-tuning. This makes it a useful starting point for larger, better-annotated, and externally validated studies.




\clearpage
\thispagestyle{plain}
\begin{center}
\noindent\normalfont\LARGE\textbf{AI Declaration}
\end{center}
\begin{quotation}
\sloppy
Generative AI tools, including OpenAI ChatGPT, were used as support tools during the work on this thesis. The tools were used for tasks such as language improvement, restructuring suggestions, LaTeX editing, consistency checks across chapters, and assistance with understanding and organizing repository content. They were also used as coding and debugging support when working with implementation material.

AI-generated suggestions were not used as a substitute for independent academic work. All methodological choices, reported results, interpretations, and final thesis text were reviewed, edited, and approved by the author. Numerical results were taken from the thesis experiments and result material. Literature-based claims were checked against the cited sources, or written cautiously when support was limited.

The author remains fully responsible for the final thesis content, including any errors, interpretations, and conclusions.
\end{quotation}


\references


\end{document}
